\documentclass[11pt]{article}
\usepackage[margin=1.1in]{geometry}
\usepackage{graphicx}
\usepackage{amsmath,amssymb,amsthm}
\usepackage{booktabs}
\usepackage{tabularx}
\usepackage[colorlinks=true,linkcolor=blue,citecolor=blue]{hyperref}
\usepackage{enumitem}

\newtheorem{proposition}{Proposition}[section]
\newtheorem{lemma}[proposition]{Lemma}
\newtheorem{remark}[proposition]{Remark}
\newtheorem{definition}[proposition]{Definition}

\newcommand{\E}{\mathbb{E}}
\newcommand{\Var}{\mathrm{Var}}
\newcommand{\Cov}{\mathrm{Cov}}
\newcommand{\lfdr}{\mathrm{lfdr}}
\newcommand{\logit}{\mathrm{logit}}
\newcommand{\one}{\mathbf{1}}

\title{Local False Discovery Rate Gated Adaptive Finite Element Regression,\\
with Moment Estimators for Undiscovered Signal:\\
\large Algorithm, Theory, and Calculations for Binomial Hazard Correction}
\author{Decision Mesh working notes}
\date{\today}

\begin{document}
\maketitle

\begin{abstract}
We describe an adaptive piecewise-linear regression procedure on a conforming triangular
mesh in which each refinement step is treated as a hypothesis and admitted through a
local false discovery rate gate. The method is developed for correcting an existing
binomial hazard model on cohort-supported state space, where most of the domain is
structurally empty and the probabilities carry unobserved pool heterogeneity. The current
development implementation separates geometric existence from statistical freedom, fits
active hierarchical coefficients by coordinate ascent on the exact tempered binomial
likelihood with a depth-specific empirical-Bayes penalty, and measures candidate evidence
by the signed root exact penalized gain. A hierarchy-aware retirement rule removes only
coefficients whose full hierarchical columns are zero on the training design. The
original empirical-null and composite-null certification results remain valid for the
legacy working-scale gate. A first attempt to transport the exact-gain score into the
same global empirical-null prefix rule failed the predeclared certified-window
instruments, despite good prediction and complete-null scale diagnostics; certification
has therefore not yet transferred to the exact solver. We record that failure and give a
conformalized held-out gain test as the next proposed certification route. The residual
moment and credibility layers quantify signal the surface has not found and decompose it
into spatially coherent and pool-specific components, with exact Gamma-frailty and
Poisson--lognormal formulas where available.
\end{abstract}

\tableofcontents

\section{Introduction}\label{sec:intro}

This paper develops an adaptive correction surface for an existing binomial hazard model. The
motivating data occupy thin cohort trajectories in a nominally two-dimensional state space, so a
standard smoother either extrapolates into empty regions or spends most of its flexibility where no
data exist. At the same time, the correction is intended for production use: every local departure
from the existing model should be supported by evidence, and the default away from observed support
should remain zero.

The procedure begins with the production-model offset, represents the correction by continuous
piecewise-linear finite elements, and proposes refinements one hierarchical coefficient at a time.
Candidate coefficients are scored, calibrated through an empirical-null two-groups model, and admitted
by a local false-discovery-rate gate. Admitted coefficients are shrunk by a depth-specific empirical
Bayes prior. A separate diagnostic layer then asks how much variation remains unexplained and divides
it into spatially coherent structure that a richer mesh might discover and pool-specific variation
that cannot be learned from the observed coordinates alone.

Figure~\ref{fig:pipeline} gives the visual map used throughout the paper. The current
development line can be remembered as
\[
\begin{aligned}
\text{tempered binomial objective}
&\longrightarrow \text{exact coordinate profiles}
\longrightarrow \text{calibrated refinement}\\
&\longrightarrow \text{hierarchically penalized fit}
\longrightarrow \text{residual attribution}.
\end{aligned}
\]
The stage fit and the candidate profile now use the same likelihood-space currency. The
remaining unresolved component is the calibration arrow: the certified results in
Sections~\ref{sec:nulls} and~\ref{sec:certnull} were obtained with the legacy
working-scale gate, and the first exact-score transfer attempt did not pass those
instruments.

The paper is organized in three passes. Section~\ref{sec:overview} gives the complete algorithm before
any derivation. Sections~\ref{sec:model}--\ref{sec:attribution} explain the statistical model,
adaptive representation, gate, shrinkage, and attribution layer. Sections~\ref{sec:scope}--
\ref{sec:realdata} then state the inferential limits and present the numerical and real-data evidence.
Technical qualifications are deliberately placed after the construction they qualify so that a first
reading can follow the main line without resolving every implementation detail.

\section{Method at a glance}\label{sec:overview}

\paragraph{Inputs.} Locations $x_i$, event equivalents $k_i$, exposures $n_i$, and an existing
log-odds offset $\eta_{0,i}$. For literal loan-count outcomes, $(k_i,n_i)$ are counts. For
balance-weighted outcomes they may be continuous variance-matching quantities.

\paragraph{Outputs.} A correction surface $\hat g$, a conforming adaptive mesh whose active
coefficients have passed the gate, estimates of residual spatial and pool variation, and optional
credibility-weighted pool effects.

\begin{enumerate}[leftmargin=2em,itemsep=2pt]
\item Rescale the coordinates, initialize a coarse conforming triangular mesh, and leave these root
scales unthresholded so later candidates have a stable neighborhood.
\item Construct the tempered binomial objective and the depth-specific hierarchical penalty. Fit the
current active coefficients by exact one-dimensional coordinate maximizations until the penalized
objective is stationary to tolerance.
\item For every admissible midpoint, build its full hierarchical column and profile the same exact
penalized objective along that coordinate. Record the optimum, the exact gain, and the signed root
gain.
\item Calibrate the candidate family and admit a multiplicity-controlled prefix. The legacy
working-scale calibration is the currently certified rule; exact-score calibration is under
revision after the failed B3 transfer experiment.
\item Activate admitted vertices, complete the newest-vertex-bisection cascades required for
conformity, and refit the active hierarchy by the exact solver.
\item Repeat profile, gate, refinement, and exact refitting until the mesh stalls across the declared
outer stages.
\item At fixed topology, retire only coefficients whose \emph{hierarchical} columns are zero on the
training design, refit once, and then estimate residual spatial and pool components and optional
credibility-weighted pool effects.
\end{enumerate}

\paragraph{Implementation status.} The exact stage solver, exact candidate profiler, and
hierarchical-support lifecycle are complete and regression tested. The signed root exact gain is
implemented in shadow mode and is the intended gate currency. Its null calibration is not yet
accepted: the first global empirical-null transfer reduced but did not eliminate false refinement on
the certified windows.

\paragraph{What is optional.} Pool attribution requires repeated observations per pool. The learned
coordinate metric, exact Gamma-frailty update, and frozen-node dynamics are extensions rather than
requirements for the core surface estimator.

\begin{figure}[ht]
\centering
\includegraphics[width=\textwidth]{fig15_pipeline.png}
\caption{The complete Decision Mesh pipeline. The working likelihood and current mesh generate candidate edge-midpoint tests; an empirical-null local-FDR gate admits supported refinements; empirical-Bayes shrinkage stabilizes the added hierarchical coefficients. After the surface stalls, a separate moment layer attributes remaining variation to missed spatial signal and pool-level components.}
\label{fig:pipeline}
\end{figure}

\section{Statistical model and effective exposure}\label{sec:model}

Data are indexed by $i=1,\dots,N$, each observation a pool-month:
covariate location $x_i\in\Omega\subset\mathbb{R}^2$ (for the motivating application
$x=(\text{WALA},\text{WAC})$ after affine rescaling of the axes to comparable ranges),
exposure $n_i$, event total $k_i$, and a known offset $\eta_{0,i}$ produced by an
existing (``internal'') model, possibly including known per-observation log-odds shifts
from observed characteristics. For literal loan-count outcomes, $n_i$ and $k_i$ are
integer counts. For balance-weighted outcomes they may instead be continuous
variance-matching pseudo-counts, as described below. The baseline generative model is
\begin{equation}\label{eq:model}
k_i \mid u_i \sim \mathrm{Binomial}\bigl(n_i,\ p_i\bigr),\qquad
\logit p_i = \eta_{0,i} + g(x_i) + u_i,
\end{equation}
where $g:\Omega\to\mathbb{R}$ is the unknown correction surface (log-odds units) and
$u_i$ is a mean-zero random effect representing unobserved variables,
$\Var(u_i)=\sigma_u^2$. The target of estimation is $g$; the targets of the diagnostic
layer (Section~\ref{sec:attribution}) are $\sigma_u^2$ and the energy of whatever part of $g$
the fit has missed.

\subsection{Continuous effective exposure for balance-weighted outcomes}
\label{sec:wacls}

The binomial representation is literal when the response is a loan count.  A
balance-weighted prepayment rate has a different sampling variance when loan
sizes differ.  Let the current balances in a pool be $b_1,\ldots,b_m$, let
$I_j\in\{0,1\}$ indicate whether loan $j$ prepays over the observation period,
and suppose first that the loans share a common conditional prepayment
probability $p$.  The prepaid balance is
\[
K_B=\sum_{j=1}^m b_j I_j,
\qquad
\Var(K_B\mid b)=p(1-p)\sum_{j=1}^m b_j^2.
\]
Define the weighted-average current loan size
\[
\mathrm{WACLS}=\frac{\sum_j b_j^2}{\sum_j b_j}.
\]
The variance-matching effective exposure and event equivalent are then
\begin{equation}\label{eq:wacls-counts}
 n_{\mathrm{eff}}
 =\frac{\sum_j b_j}{\mathrm{WACLS}}
 =\frac{(\sum_j b_j)^2}{\sum_j b_j^2},
 \qquad
 k_{\mathrm{eff}}
 =\frac{\sum_j b_j I_j}{\mathrm{WACLS}}.
\end{equation}
They satisfy $\E[k_{\mathrm{eff}}\mid b]=n_{\mathrm{eff}}p$ and
$\Var(k_{\mathrm{eff}}\mid b)=n_{\mathrm{eff}}p(1-p)$ exactly.  Thus
$(k_{\mathrm{eff}},n_{\mathrm{eff}})$ has the first two moments of a binomial
observation even though neither quantity need be an integer.  The fitting
criterion
\[
\ell(\eta)=k_{\mathrm{eff}}\eta-n_{\mathrm{eff}}\log(1+e^\eta)
\]
is therefore used as a quasi-binomial log likelihood.  All score, IRLS, and
Pearson-residual formulas below remain unchanged after replacing $(k_i,n_i)$
by continuous $(k_{\mathrm{eff},i},n_{\mathrm{eff},i})$.

The effective count is the inverse Herfindahl concentration of current loan
balances.  Equal-sized loans give $n_{\mathrm{eff}}=m$; concentration in a few
large loans reduces the information.  For loan-size coefficient of variation
$c_v$ and a moderately large pool,
$n_{\mathrm{eff}}/m\approx(1+c_v^2)^{-1}$.  Dividing prepaid balance by the
ordinary average loan size while retaining the literal loan count therefore
understates the variance by approximately $1+c_v^2$.  Equation
\eqref{eq:wacls-counts} is intended for balance-weighted outcomes only; when
actual loan counts are observed, the disclosed count remains preferable.

\begin{figure}[ht]
\centering
\includegraphics[width=0.96\textwidth]{fig19_wacls_effective_count.png}
\caption{Why balance-weighted outcomes require an effective exposure. Equal-sized loans preserve the literal count, while concentration of the same total balance in a few large loans reduces the variance-matching exposure. The WACLS formula is the inverse concentration of current balances and exactly matches the first two moments when loans share a conditional event probability.}
\label{fig:wacls-concept}
\end{figure}

Two features of the application drive the design. First, the support of the sampling
distribution is a union of cohort trajectories: at any calendar month each loan age is
occupied by exactly one vintage, so cross-sections are one-dimensional curves and the
two-dimensional support is swept out over time. Large regions of $\Omega$ are
structurally empty and must never receive confident extrapolation. Second, the
correction is a \emph{deliverable against a production model}: every basis function that
exists in $\hat g$ should be individually certified, and the default in the absence of
evidence must be $\hat g=0$.


\begin{figure}[h]
\centering
\includegraphics[width=\textwidth]{fig2_data.png}
\caption{The simulated data regime. Left: observation density in $(\text{WALA},\text{WAC})$ under cohort dynamics; support is confined to diagonal cohort trajectories and dies with attrition. Right: fitted $\hat p$ spans essentially $(0,1)$ (the wide-$p$ regime), which is what makes the re-expansion of Section~\ref{sec:glm} and the $(1-2\hat p)$ terms of Section~\ref{sec:attribution} load-bearing.}
\label{fig:data}
\end{figure}

\noindent\textbf{Assumption (no common calendar shocks).} The model contains
no calendar-time effects: conditional on $(\text{WALA},\text{WAC})$ and the
pool effects, observations in the same calendar month are independent. Real
prepayment data violates this; rate-path surprises, seasonality, and media
effects are common monthly shocks, neither i.i.d.\ pool noise nor functions
of the coordinates; and because pools in a cell are largely observed in
the same months (vintage and coupon are correlated), such shocks load onto
the cross-pool moments and, if serially correlated, contaminate the pool
moments. This assumption is stated as a limitation: the simulator cannot
reveal the violation on its own because its generator matches the model, and
calendar-month intercepts in the fit (standard in production prepayment
models) are the appropriate extension when multi-month real panels are in
scope.

\paragraph{Status and reproducibility.} Simulation claims in this document
are labeled with their seed counts; unlabeled figures are single-seed
illustrations, and every estimator-level claim rests on three to ten seeds
as stated at the point of use. Section~\ref{sec:conditionality} tabulates
which inferential outputs are conditional on which estimated quantities;
claims should be read against that table.

\section{Likelihood and one-step scoring}\label{sec:glm}

\subsection{Offset GLM and IRLS}
Conditional on $u\equiv 0$, \eqref{eq:model} is a binomial GLM with offset $\eta_0$ and
linear predictor $\eta_i=\eta_{0,i}+\sum_v h_v\varphi_v(x_i)$. Fisher scoring for the
log-likelihood $\ell(h)=\sum_i k_i\eta_i-n_i\log(1+e^{\eta_i})$ is iteratively
reweighted least squares on the working response and weights
\begin{equation}\label{eq:working}
y_i \;=\; \hat g^{(m)}(x_i)+\frac{k_i-n_i p_i^{(m)}}{n_i p_i^{(m)}q_i^{(m)}},
\qquad
w_i \;=\; n_i p_i^{(m)} q_i^{(m)},
\qquad p_i^{(m)}=\sigma\!\bigl(\eta_{0,i}+\hat g^{(m)}(x_i)\bigr),
\end{equation}
with $q=1-p$ and $\sigma$ the logistic function. The inner problem at each outer
iteration $m$ is weighted least squares in the P1 basis; at a fixed point the heights
satisfy the binomial score equations, so the estimator is maximum (penalized) likelihood
implemented on a WLS engine. In practice one outer step at the internal model
($m=0$, $\hat g^{(0)}=0$) captures most of the correction and outer refreshes are
applied whenever the refinement gate stalls.

\begin{proposition}[first-order unbiasedness of one-step scoring]\label{prop:onestep}
Let $\logit p=\eta_0+\delta$ and $y=(k-np_0)/(np_0q_0)$ with $p_0=\sigma(\eta_0)$. Then
$\E[y]=\delta+O(\delta^2)$ and $\Var(y)=1/(np_0q_0)+O(\delta)$.
\end{proposition}

\emph{Adequacy criterion.} The $O(\delta^2)$ bias coefficient
$\tfrac12(1-2p_0)$ is $\approx\tfrac12$ at hazard scale, so a correction of
$0.8$ log-odds carries a first-round bias near $0.3$: one-step is a
coarse-regime device. The operational criterion used throughout is: one-step
is adequate while the largest admitted surplus per round stays below
$\sim0.5$ log-odds; beyond that, the IRLS re-expansion (every third round and
on stall) re-linearizes, and the wide-$p$ regimes in this document always run
with re-expansion on. No claim is made for one-step alone outside that
envelope.
\begin{proof}
$\E k=np(\delta)$ and $p(\delta)=p_0+p_0q_0\delta+O(\delta^2)$ by
$dp/d\eta=pq$; divide by the information $np_0q_0$. The variance statement follows from
$\Var k=np(\delta)q(\delta)$.
\end{proof}

\subsection{Exact tempered penalized coordinate ascent}\label{sec:exactsolver}
The IRLS construction above remains useful for diagnostics and for the legacy gate, but it is no
longer the stage height solver on the development line. Let $c_i\in(0,1]$ be the current
pool-law tempering factor and let $\lambda_{d(v)}$ be the precision of the depth-specific Gaussian
hierarchical prior. With prior targets $m_v$ held fixed during a sweep, the stage objective is
\begin{equation}\label{eq:exact-objective}
Q(h)=\sum_i c_i\left[k_i\eta_i-n_i\log(1+e^{\eta_i})\right]
-\frac12\sum_{v\in V_{\mathrm{free}}}\lambda_{d(v)}(\delta_v-m_v)^2,
\qquad
\eta_i=\eta_{0,i}+\hat g_h(x_i).
\end{equation}
For coordinate $v$, the solver maximizes the concave one-dimensional function
$Q_v(a)=Q(h+a e_v)$ by safeguarded Newton steps and accepts an update only when the full objective
is nondecreasing. The same hierarchical column engine is used by the stage solver, candidate
profiler, information census, and retirement pass.

This change removes the eruption mechanism rather than merely clipping it. A saturated binomial
likelihood may have tiny local Fisher curvature, so a stale quadratic can propose an enormous move;
the exact objective still has a finite penalized optimum and a finite gain. In the B1 closure tests,
thousands of committed coordinates were reprofiled against the full objective with zero failures,
and the stored gain and realized objective increase agreed to roughly $10^{-10}$ nats. After removal
of the now-dead saturation fallback, a 20-split battery had zero eruptions, mean held-out NLL
$2.9449$, standard deviation $0.0054$, and worst split $2.9537$.

\subsection{Why the empirical logit must not be used}\label{sec:empiricallogit}
The classical alternative response, the Haldane--Anscombe empirical logit
$\tilde y=\log\{(k+\tfrac12)/(n-k+\tfrac12)\}$ with Gart variance
$(k+\tfrac12)^{-1}+(n-k+\tfrac12)^{-1}$, is designed for $npq$ moderately large. At
prepayment scale ($p\sim 0.2$--$2\%$ monthly, $n$ of order $10^2$) it fails for two
compounding reasons. (i) \emph{Floor}: $\tilde y\ge \log\{0.5/(n+0.5)\}$
($=-6.22$ at $n=250$) regardless of how small $p$ is, and the flooring event $k=0$ is
the modal outcome; e.g.\ at $np=1$, $P(k{=}0)\approx e^{-1}\approx 0.37$. (ii)
\emph{Outcome-dependent weights}: the Gart precision is approximately $k+\tfrac12$, so
observations that randomly realized more events receive both larger $\tilde y$ and more
weight, biasing any precision-weighted average upward. In the development simulation the
combined bias of the precision-weighted mean was $\approx +0.5$ log-odds against a truth
of $-0.82$; the mesh reproduced its (biased) target almost exactly, capping recoverable
signal at $59\%$, versus $93$--$94\%$ under one-step scoring. The one-step construction
\eqref{eq:working} has no floor (a $k=0$ month contributes a mildly negative residual)
and its weights are functions of the model, not of $k$.

\begin{remark}[the $+\tfrac12$ in its proper place]
Regularization against separation is still required at the \emph{likelihood} level: a
vertex whose entire star contains only $k=0$ observations has a monotone likelihood and
unpenalized ML sends $h_v\to-\infty$, dragging every adjacent face with it because the
basis is continuous. Firth's bias-reduction penalty (the Jeffreys prior) is the
principled fix \cite{firth1993}; on a single binomial cell it reduces exactly to adding
$\tfrac12$ to each count. Applied once, at the likelihood, on pooled counts, the idea is
correct; applied per observation to a plug-in estimator with outcome-dependent weights,
it is the failure described above. In the implementation the empirical Bayes prior of
Section~\ref{sec:eb} supplies the penalty.
\end{remark}

\section{Overdispersion and known covariates}\label{sec:overdisp}

\subsection{Random effects in $p$}
With $u_i$ present, the working residual has variance $1/w_i+\sigma_u^2$, so the
efficient weights are the DerSimonian--Laird weights
\begin{equation}\label{eq:dlw}
\tilde w_i=\Bigl(\frac{1}{w_i}+\sigma_u^2\Bigr)^{-1}\ \le\ \frac{1}{\sigma_u^2}.
\end{equation}
The saturation bound is operationally decisive for MBS data, where $n$ spans single
digits (small spec pools) to $10^5$ (Ginnie~II multi-issuer): a mega-pool with
$w=npq\approx 10^3$ carries $\tilde w\approx 1/\sigma_u^2\approx 11$ at
$\sigma_u=0.3$, i.e.\ is worth only a few hundred-loan pools; cross-sectional pool
count, not loan count, drives precision. Uncorrected, mega-pool null $z$'s are inflated
by $\sqrt{1+w\sigma_u^2}$, far beyond the empirical null's constraint, and the gate
over-refines into $u$-noise. At small $w\sigma_u^2$ (the homogeneous-$n$ simulations,
inflation $\approx 11\%$) the constrained empirical null absorbs the effect and the
correction is nearly redundant; heterogeneity in $w$ is what makes it load-bearing.

\subsection{Known log-odds shifts}
A characteristic $c_i$ with \emph{known} coefficient enters the offset,
$\eta_{0,i}\mapsto\eta_{0,i}+\beta c_i$; the entire pipeline is invariant to it (weights
become heteroskedastic per observation, which the weighted face statistics carry
through). If instead the shift is omitted, it is independent of location and therefore
statistically indistinguishable from $u$:

\begin{proposition}[misspecification identity]\label{prop:audit}
If a known shift $\beta c_i$, $c_i\perp x_i$, $\Var(c)=1$, is omitted from the offset,
the idiosyncratic moment of Section~\ref{sec:attribution} estimates
$\sigma_u^2+\beta^2$ in place of $\sigma_u^2$.
\end{proposition}

\noindent This yields two tools: an \emph{audit} (measured $\sigma_u^2$ above what is
plausible for genuine unobservables quantifies characteristics missing from the
offset) and a \emph{valuation statistic} (the drop in $\hat\sigma_u^2$ when a candidate
adjustment is added to the offset is the log-odds variance that characteristic
explains). Development numbers: $\hat\sigma_u^2=0.082$ with the shift in the offset
versus $0.201$ without, against $0.090$ and $0.090+0.35^2=0.2125$.

\section{Adaptive finite-element representation}\label{sec:representation}

\subsection{Newest vertex bisection}
$\Omega$ is triangulated starting from two right triangles obtained by splitting the
bounding box along a diagonal. Refinement is by newest vertex bisection (NVB): each
triangle carries a refinement edge (the edge opposite its newest vertex); activating the
midpoint $v$ of an edge $e$ bisects $e$ and both adjacent triangles along the chord from
$v$ to each opposing vertex. If a triangle is asked to split along a non-refinement
edge, a completion cascade first bisects it along its own refinement edge, recursively;
the mesh therefore remains conforming at all times \cite{mitchell1991,binev2004}. Under
NVB all triangles produced from a right isosceles initial mesh belong to a finite set of
similarity classes, so shape regularity (minimum angle bounded away from zero) holds
uniformly at every depth; this is why the axes must be pre-scaled: on a raw domain of
aspect ratio $\sim 14$, every admissible bisection violates the shape guard and
refinement cannot start.

Point location is by a binary tree recording, at each internal node, the separating
line between its two children. The separating line of an NVB split is the \emph{chord}
(midpoint to opposing vertex), not the bisected edge; the parent face lies entirely on
one side of the line containing its own edge, so storing the edge yields a walk that is
deterministic and wrong at every non-root node. With the chord stored, prediction is an
$O(\log F)$ walk followed by barycentric interpolation.

\subsection{Basis and surplus coefficients}
Write $V$ for the active vertices and $\{\varphi_v\}_{v\in V}$ for the continuous
piecewise-linear (P1) hat functions on the current mesh, $\varphi_v(v')=\delta_{vv'}$.
The fitted surface is $\hat g(x)=\sum_{v} h_v\varphi_v(x)$. For a midpoint $v$ of edge
$e=(p_0,p_1)$, define the \emph{linear interpolant value}
$\mu_{\mathrm{lin}}(v)=\tfrac12(h_{p_0}+h_{p_1})$ and the \emph{surplus} (hierarchical
coefficient)
\begin{equation}\label{eq:surplus}
\delta_v \;=\; h_v-\mu_{\mathrm{lin}}(v).
\end{equation}
Activating $v$ adds exactly one degree of freedom, and $\delta_v=0$ reproduces the
pre-split surface; the hierarchy $\{\delta_v\}$ plays the role of wavelet coefficients
in \cite{abramovich2006}, which is the closest antecedent of the gating idea below.

\begin{figure}[ht]
\centering
\includegraphics[width=\textwidth]{fig16_hierarchical_surplus.png}
\caption{The hierarchical coefficient tested by a candidate edge split. Before refinement the midpoint value is fixed by interpolation between its parents. Activating the midpoint adds one local degree of freedom, the surplus $\delta_v=h_v-\mu_{\mathrm{lin}}(v)$; setting $\delta_v=0$ exactly reproduces the pre-split surface.}
\label{fig:surplus-concept}
\end{figure}

\subsection{Geometry, statistical freedom, and hierarchical support}\label{sec:lifecycle}
A vertex can exist geometrically without carrying an independent coefficient. This distinction is
necessary because NVB conformity cascades create midpoint vertices that are required by the mesh but
were not themselves admitted by the gate. The implementation therefore distinguishes:
\begin{enumerate}[leftmargin=2em,itemsep=1pt]
\item \emph{geometry-only or constrained}: the vertex exists and its surplus is fixed at zero;
\item \emph{free}: the surplus is an optimized statistical coordinate;
\item \emph{dormant}: a free coordinate currently has negligible likelihood information but is not
structurally absent;
\item \emph{retired}: at fixed topology the coordinate has been proved to have zero training support
and is returned to zero surplus.
\end{enumerate}

For a free coefficient $v$, define its hierarchical design column
\begin{equation}\label{eq:hiercol}
X^H_{iv}=\frac{\partial \hat g(x_i)}{\partial \delta_v}.
\end{equation}
The derivative propagates through constrained descendants and stops when it reaches another free
coefficient. It is therefore generally different from the nodal P1 hat function. Structural support,
likelihood information, and replication are separate quantities:
\begin{equation}\label{eq:support-info}
S_v^{\mathrm{geom}}=\sum_{i\in\mathcal T}(X^H_{iv})^2,
\qquad
I_v=\sum_{i\in\mathcal T} c_i n_i p_i(1-p_i)(X^H_{iv})^2,
\end{equation}
where $\mathcal T$ is the explicit training split. Retirement uses only
$S_v^{\mathrm{geom}}$: small $I_v$ makes a coordinate imprecise, and low effective pool count makes
it difficult to certify, but neither makes the coordinate geometrically nonexistent.

Retirement is performed deepest first, rebuilding the hierarchy after each depth batch. Before any
post-retirement refit, each batch must leave all training predictions invariant to numerical
precision. On the original-split B1 regression this rule retired $67$ genuinely empty columns,
rescued $12$ vertices that were nodally empty but hierarchically supported, and changed the training
predictor by exactly zero before reoptimization. The previous nodal-star cleanup had incorrectly
retired supported ancestors and moved some predictions by several logits.

\begin{figure}[ht]
\centering
\includegraphics[width=\textwidth]{fig20_mesh_build_progression.png}
\caption{Mesh construction across the three outer fitting stages for one representative continuous-WACLS run. The first panel is the common pre-refined mesh. Colored points in later panels mark vertices added since the preceding panel. Refinement remains concentrated on observed support, while variance reweighting between stages changes which local deviations remain worth resolving.}
\label{fig:mesh-build}
\end{figure}

\section{Local-FDR-gated refinement}\label{sec:gate}

Each inactive midpoint $v$ (a candidate refinement) is a hypothesis
\[
H_0(v):\ \delta_v=0 \quad\text{(the surface is linear across edge $e(v)$)} .
\]
Fixing all other heights, the coordinate-wise WLS estimate of $h_v$ over the (simulated)
star of $v$ is
\begin{equation}\label{eq:coord}
\hat\beta_v=\frac{x^\top W r}{x^\top W x},\qquad
x=(\varphi_v(x_i))_i,\ \ W=\mathrm{diag}(w_i),\ \
r_i=y_i-\!\!\sum_{u\neq v}\!h_u\varphi_u(x_i),
\end{equation}
with conditional variance $\sigma^2/(x^\top Wx)$, where $\sigma^2$ is the working-scale
noise ($\sigma=1$ when $w$ carries the full precision, as in \eqref{eq:working}).

\subsection{Exact candidate gain and its signed-root coordinate}\label{sec:exactgain}
For an inactive candidate, let $X^H_v$ be the hierarchical column it would acquire upon activation
and profile the exact objective \eqref{eq:exact-objective} while holding the existing coordinates
fixed:
\begin{equation}\label{eq:exactgain}
\hat a_v=\arg\max_a Q_v(a),
\qquad
\Delta Q_v=Q_v(\hat a_v)-Q_v(0),
\qquad
r_v=\operatorname{sign}(\hat a_v)\sqrt{2\Delta Q_v}.
\end{equation}
Because zero is feasible, $\Delta Q_v\ge0$ up to numerical tolerance. If the objective is locally
quadratic, $Q_v(a)\approx Q_v(0)+s_va-\tfrac12J_va^2$, then
$r_v=s_v/\sqrt{J_v}$: the signed root gain is exactly the nonlinear continuation of the familiar
score/Wald statistic. Away from the trust region it follows the actual objective rather than the
local tangent parabola.

The penalty is included deliberately. Thus $r_v$ is a direction-preserving measure of net
penalized objective improvement, not automatically a classical $N(0,1)$ likelihood-ratio pivot. Its
null law must be calibrated under the same tempering, estimated pool law, adaptive fit state, and
candidate geometry. B2 shadow tests found the implementation internally exact: across $24{,}404$
original-split profiles there were no negative gains or sign inconsistencies, and
$r_v^2=2\Delta Q_v$ held to approximately $6\times10^{-11}$. The new and legacy currencies are not
interchangeable: their signed correlation was $0.253$, $31\%$ of candidate directions disagreed,
and only about $31\%$ of their top deciles overlapped within a round. At the named regime-straddle
candidate near raw $(0,4.52)$, however, the two agreed: the exact gain was $4.829$ nats and
$r_v=-3.108$ versus legacy $z=-3.059$.

The following proposition describes the \emph{legacy working-scale score}. It remains relevant
because all currently accepted certified-null counts were produced by that gate, and because its
variance model supplies one candidate-specific scale estimate for exact-score calibration.

\begin{proposition}[null distribution with estimated neighbors]\label{prop:null}
Suppose on $\mathrm{supp}\,\varphi_v$ the true surface equals the current interpolant
except for independent estimation errors in the parent heights,
$\Var(\hat h_{p_j})=\sigma^2 s_{p_j}$, approximately independent of the data in the
star. Then under $H_0(v)$,
\begin{equation}\label{eq:z}
z_v=\frac{\hat\beta_v-\hat\mu_{\mathrm{lin}}(v)}
{\sigma\sqrt{\dfrac{1}{x^\top Wx}+\dfrac{s_{p_0}+s_{p_1}}{4}}}
\;\xrightarrow{\ d\ }\;N(0,1).
\end{equation}
\end{proposition}

\emph{Caveat (neglected cross-covariance).} The independence approximation
neglects the covariance between the parent estimates and the candidate
statistic; since a candidate's star is contained in the union of its parents'
stars, the parents' estimation error is built partly from the candidate's own
data, and the same estimates enter both $\hat\mu_{\mathrm{lin}}$ and, through
the residuals, $\hat\beta_v$. The neglected term has definite sign structure
and grows with the star's share of its parents' information; negligible at
coarse depths, not at fine depths near the information floor. A bound in
terms of that information share is the missing result; until it exists, the
proposition should be read as a calibration approximation whose quality
degrades with depth, one function the empirical-null re-centering (central
matching) partially absorbs. The bound now exists, with a verification that
reorders its importance.

\emph{Proposition (covariance bound, local-WLS idealization).} With
$\hat\beta_v=\sum w\varphi_v y/I_v$, $\hat h_p=\sum w\varphi_p y/I_p$, and
$J_{vp}=\sum_i w_i\varphi_v\varphi_p\ge0$, the neglected covariance is
exactly $\mathrm{Cov}(\hat\beta_v,\hat h_p)=J_{vp}/(I_vI_p)\ge0$, and by
Cauchy--Schwarz over $\mathrm{star}(v)$,
\[
\mathrm{Corr}(\hat\beta_v,\hat h_p)\;\le\;\sqrt{\rho_p},\qquad
\rho_p=\frac{\text{parent $p$'s information inside }\mathrm{star}(v)}{I_p},
\]
tight when $\varphi_p\propto\varphi_v$ on the star; chaining gives
$|V/\hat V-1|\le3\sqrt{\rho^*}$ with $\rho^*$ the largest of the two
candidate--parent shares and the parent--parent share. Verified against
Monte Carlo on heterogeneous-weight geometries to three--four decimals; the
candidate--parent (deflation) channel dominates the parent--parent
(inflation) channel, with $V/\hat V\in[0.81,0.97]$ at shares up to $0.22$
--- so within the idealization the neglect makes the test mildly
\emph{conservative}, degrading as $\sqrt{\rho^*}$ with depth, as claimed.

\emph{In-vivo refutation of the idealization's relevance.} The idealization
predicts null scores mildly under-dispersed and worsening with depth.
Measured on pure-null runs, the candidate scores are instead
\emph{over}-dispersed pooled ($\mathrm{sd}\approx2.0$) and flat in depth ---
because the pooled figure mixes rounds whose null scale ranges from
$0.39$ to $3.6$: scores blow up in the rounds immediately following stage
transitions (working-weight changes leave the residual field out of
equilibrium) and collapse immediately after re-expansion sweeps. Null
calibration is therefore dominated by \emph{fit-state transients}, not by
the neglected covariance; the $\sigma_0\le1.4$ box binds in the transition
rounds, which is the concrete candidate mechanism for the complete-null FDP
excursions measured above. The covariance bound is thus established but
third-order in practice; the first-order repair is procedural; score the
gate only at sweep-converged states, or widen the empirical-null box in
transition rounds; and is recorded as the queued engineering item.
\begin{proof}[Calculation]
$\hat\mu_{\mathrm{lin}}=\tfrac12(\hat h_{p_0}+\hat h_{p_1})$ has variance
$\tfrac14\sigma^2(s_{p_0}+s_{p_1})$ under the independence approximation, and
$\Var(\hat\beta_v)=\sigma^2/(x^\top Wx)$; the cross-covariance is neglected (the parent
estimates are dominated by data outside the candidate's star). Normality is the usual
weighted CLT for \eqref{eq:coord}. The correction matters: omitting it, null $z$'s are
overdispersed by the estimated-neighbor variance and, in the development experiments,
the empirical null (Section~\ref{sec:lfdr}) then either swallowed genuine signal or,
with a theoretical $N(0,1)$ null, admitted essentially all candidates and diverged.
\end{proof}

Two practical guards accompany the test. A candidate is tested only if its star
contains at least a minimum number of observations (information floor), and all $z$'s
are winsorized at $|z|=9$ before density estimation (any such candidate is a certain
discovery; winsorization only protects the histogram).

\subsection{Empirical-null calibration and local-FDR admission}\label{sec:lfdr}

\subsection{Two-groups model and Lindsey estimation}
At each refinement round, collect the candidate scores $z_1,\dots,z_M$ and posit the
two-groups model \cite{efron2004,efron2008}
\[
f(z)=\pi_0 f_0(z)+(1-\pi_0)f_1(z),\qquad
\lfdr(z)=\frac{\pi_0 f_0(z)}{f(z)}=P(H_0\mid z).
\]
The marginal $f$ is estimated by Lindsey's method: histogram counts $c_b$ on bins with
midpoints $m_b$ are modeled as independent Poisson with
$\log \E c_b=\sum_{k=0}^{6}\beta_k\,\psi_k(m_b)$ (standardized polynomial basis), fit by
IRLS for the Poisson GLM; $\hat f$ is the normalized exponentiated fit.

\subsection{Empirical null by central matching}
Because \eqref{eq:z} is approximate (neighbor covariances, EB shrinkage, unconverged
heights), the null is estimated rather than assumed. If near the center
$f(z)\approx \pi_0\,\phi\!\bigl((z-\delta_0)/\sigma_0\bigr)/\sigma_0$, then
$\log f(z)=a_0+a_1z+a_2z^2$ with
\begin{equation}\label{eq:centralmatch}
\sigma_0^2=-\frac{1}{2a_2},\qquad
\delta_0=a_1\sigma_0^2,\qquad
\pi_0=\exp\!\Bigl(a_0+\frac{\delta_0^2}{2\sigma_0^2}\Bigr)\sqrt{2\pi}\,\sigma_0 ,
\end{equation}
so a quadratic fit of $\log\hat f$ on $|z|\le 2$ identifies the null. Central matching
assumes the center is null-dominated; mid-refinement, when most candidates are true
signal, an unconstrained fit reads the wide center as a wide null and halts refinement
prematurely (observed directly in development). Since the variance correction in
\eqref{eq:z} already removes the first-order inflation, the empirical null is
\emph{constrained} to mop up residual dispersion only:
$\sigma_0\in[1,1.4]$, $\delta_0\in[-0.5,0.5]$, $\pi_0\in[0.5,1]$.

\subsection{Selection rule and FDR}
Sort candidates by $\widehat{\lfdr}$ and admit the largest prefix $S$ with
$\frac{1}{|S|}\sum_{i\in S}\widehat{\lfdr}(z_i)\le q$ (target $q=0.10$).

\begin{figure}[ht]
\centering
\includegraphics[width=0.98\textwidth]{fig17_lfdr_gate.png}
\caption{A representative local-FDR refinement round. Candidate scores are calibrated against an empirical null; low-lfdr tail candidates are ranked, and the longest prefix whose mean estimated local FDR stays below $q$ is admitted. The figure explains the operational rule; formal qualifications for dependence and adaptive hypothesis generation appear below.}
\label{fig:lfdr-concept}
\end{figure}

\begin{proposition}[Bayes FDR of the prefix rule]\label{prop:fdr}
Under the two-groups model with known $(\pi_0,f_0,f)$,
$\E[\mathrm{FDP}(S)\mid z_{1:M}]=\frac{1}{|S|}\sum_{i\in S}\lfdr(z_i)\le q$ for the
selected set, i.e.\ the rule controls the posterior expected false discovery
proportion at $q$.
\end{proposition}

\emph{Scope of Proposition~\ref{prop:fdr} (downgrade).} Three qualifications
bound what this controls, per an external review we credit. (i) The candidate
$z$'s within a round are positively dependent (adjacent stars share parents
and observations); the expectation argument for Bayes FDR survives dependence,
but the Lindsey estimation of $\hat f$ does not cleanly; histogram counts
are no longer independent Poisson, $\hat f$ has variance the machinery does
not see, and the realized FDP can be substantially more variable than its
expectation. (ii) The floor $\pi_0\ge0.5$ is deliberately conservative and is
wrong mid-refinement when most candidates are true signal, so ``mean lfdr
$\le q$'' then controls a quantity more conservative than the nominal one.
(iii) Most fundamentally, the hypotheses are random: $H_0(v)$ is a functional
of the current mesh, which is a functional of the data through prior rounds
--- a hypothesis-set selection problem outside the fixed-dictionary setting of
Abramovich et al. The honest statement of the proposition is therefore:
\emph{per-round Bayes FDR under the estimated lfdr, conditional on the mesh at
that round}, with whole-run frequentist FDP control an empirical matter, now measured.

\emph{Empirical FDP study (10 seeds per cell).} A gate admission is classified
false when the true hierarchical surplus at the vertex,
$|f(v)-\tfrac12(f(p_0)+f(p_1))|$, is below $\varepsilon$ (primary
$\varepsilon=0.01$ log-odds; results bracketed at $0.005/0.02$);
pre-refinement vertices are excluded, only lfdr-gate admissions count. In
signal regimes the realized FDP is conservative and nearly flat in $q$:
means $0.060/0.062/0.066$ at $q=0.05/0.10/0.25$ on the uniform surface
(admission counts $293/331/465$), $0.068$ under pure white pool noise and
$0.056$ on the chirp at $q=0.10$, with per-seed maxima $\le0.134$; the
conservatism stack ($\pi_0$ floor, dispersion calibration, empirical null)
dominates the nominal target, and at $q=0.05$ the mean mildly exceeds nominal
($0.060$), so control should not be described as tight at the small-$q$ end.
At the \emph{complete null} (zero surface), however, the per-round Bayes
statement visibly fails to deliver whole-run FDP control: seven to eight of
ten seeds admit at least one vertex ($36\pm62$ admissions with frailty
present, $71\pm116$ without, out of a candidate stream in the tens of
thousands), and any admission at the null makes $\mathrm{FDP}=1$, so the
realized FDP is approximately Bernoulli with mean $0.7$--$0.8$ against
$q=0.10$. The absolute false counts are small and the admitted heights tiny
(the no-signal energy floor is unaffected), but the qualitative point stands:
with an estimated null and dependent scores, FDP at or near the complete null
is an excursion event, not a controlled rate; exactly the fragility the
dependence caveat above predicts. Downstream consumers should read the gate
as a conservative discovery filter in signal regimes and as a
small-but-uncontrolled admitter at the null.
\begin{proof}
$\E[\#\{\text{true nulls in }S\}\mid z_{1:M}]=\sum_{i\in S}P(H_{0,i}\mid z_i)
=\sum_{i\in S}\lfdr(z_i)$ by definition of lfdr and (conditional) independence; divide
by $|S|$.
\end{proof}

\begin{remark}[caveats]
Control is exact for the oracle lfdr and one exchangeable batch; in the algorithm
(i) $\lfdr$ is estimated, (ii) rounds are sequential and each round's candidate set
depends on past admissions. The constrained empirical null is conservative for (i);
for (ii) the batch-round structure (test all current candidates, admit, re-converge,
re-test) is the finite-sample analogue of the wavelet-coefficient setting of
\cite{abramovich2006}, and online-FDR procedures (alpha-investing, SAFFRON) are the
fully rigorous alternative if formal sequential control is required.
\end{remark}

\subsection{Coarse scales are not thresholded}
The coarsest hat coefficients of an oscillatory target are near zero (the signal is
nearly orthogonal to the coarse basis), so a gate applied from depth zero can fail to
start: with few candidates and small coarse coefficients, nothing is admitted and the
procedure stops at the initial mesh (observed: 2 faces, entire signal missed).
Following the wavelet-thresholding convention, the mesh is uniformly pre-refined to a
modest face count ($\approx 48$) before gating; the cost is a handful of possibly-null
degrees of freedom, absorbed by the EB prior, against removal of the cold-start failure
mode. For small candidate sets ($M<50$) the Lindsey fit is unstable and a
Benjamini--Hochberg step at level $q$ on two-sided $p$-values replaces it
\cite{bh1995}.

\subsection{Round structure and stopping}
Each round: (1) re-estimate the EB hyperparameters; (2) run coordinate-descent height
sweeps to a tolerance so that candidate residuals are computed against converged
neighbors (unconverged heights contaminate every $z$ in a star); (3) score all
candidates, fit the two-groups model, admit the prefix set; (4) activate admitted
vertices (NVB cascades may add completion vertices, which are structural, not
discoveries). If nothing is admitted, deepen the height convergence and re-test; if the
binomial outer loop is in use, re-expand \eqref{eq:working} at the current fit and
re-test; stop after repeated stalls. The stopping rule is thus intrinsic: refinement
ends when no remaining candidate clears the FDR gate, not when a time or size budget is
exhausted.

\section{Empirical Bayes fitting and shrinkage}\label{sec:eb}

\noindent\emph{Selection bias in $\hat\tau_d^2$ (acknowledged, uncorrected).}
The depth-level moment for $\tau_d^2$ is computed over \emph{admitted}
vertices, whose $\hat\delta_v$ are truncated Gaussians (selected for $|z|$
beyond a data-dependent threshold), not the unshifted Gaussians the moment
identity assumes. Truncation inflates the moment mechanically; the
selected-null contribution alone is $\mathbb{E}[z^2\,|\,\text{selected}]
\approx4$--$9$ at $q=0.10$; so $\hat\tau_d^2$ is biased up by an amount
depending on the threshold and the true slab; the same effect corrected in
the variance-component moments is presently uncorrected here. Two further
honest notes: shrinkage under admitted structure leaves signal-adjacent
residual that biases neighboring candidates' scores in later rounds, a
feedback the theory does not treat; and dropping the $\sum w^2/\sum w$ term
shrinks $\hat\tau^2$, which is conservative against false structure but
anti-conservative for the missed-signal moments downstream, which read the
extra shrinkage as undiscovered energy. The truncated-moment correction is now derived,
verified in isolation, and; instructively; found not to help in the
live loop. For a centered Gaussian $X\sim\mathcal N(0,V)$ selected on
$|X|>a$, $\mathbb E[X^2\mid|X|>a]=V\kappa(a/\sqrt V)$ with
$\kappa(\alpha)=1+\alpha\varphi(\alpha)/(1-\Phi(\alpha))$, so the
selection-consistent depth moment solves
$\chi^2_d=\sum_v\kappa_v(\tau^2)\,(\tau^2+s_v)/s_v-1$ for $\tau^2$
(bisection; reduces to plain DL at $a_v=0$), with $a_v$ recorded at
admission as $c_{\mathrm{round}}$ times the candidate's calibrated standard
error. On synthetic draws with heterogeneous $s$ the naive moment is biased
upward by $+6$ to $+17$ variance units at realistic thresholds; the
mechanism quantified above; while the corrected estimator lands within
$\pm0.09$ of truth, including under $10\%$ null contamination. Deployed
inside the pipeline, however, the correction \emph{worsens} held-out
deviance on two of three seeds and destabilizes the gate on a third, and is
neutral on real data. The cause is localized: the pipeline's moment consumes
already-shrunk (pooled) coefficients, whose $\chi^2$ is deflated relative to
the raw $\hat\delta$ the truncation law describes, so the naive moment's
selection inflation and the shrunk-input deflation partially cancel;
correcting only the former over-corrects, and the smaller $\hat\tau^2$
feeds harder shrinkage and gate feedback. The estimator therefore ships
opt-in (\texttt{DMESH\_TRUNC=1}), the honest theoretical status is that the
prior-estimation moment carries two biases of opposite sign whose partial
cancellation is currently load-bearing, and the principled repair is the
truncated moment applied to \emph{raw admission surpluses}; queued, as it
requires carrying raw coefficients alongside pooled ones.

Admitted surpluses are shrunk under a Gaussian working prior per depth $d$:
$\delta_v\sim N(0,\tau_d^2)$, giving the posterior-mean update
\begin{equation}
\hat h_v=\mu_{\mathrm{lin}}(v)+\frac{\tau_d^2}{\tau_d^2+\sigma^2 s_v}\,
\bigl(\hat\beta_v-\mu_{\mathrm{lin}}(v)\bigr),\qquad s_v=\frac{1}{x^\top Wx}.
\end{equation}
Hyperparameters are estimated by the DerSimonian--Laird moment identity
\cite{dl1986}: for $m$ vertices at depth $d$ with weights $w_v=1/s_v$ and
$Q=\sum_v w_v(\hat\delta_v-\bar\delta_w)^2$,
\begin{equation}\label{eq:dl}
\E Q=(m-1)\sigma^2+\tau_d^2\Bigl(\textstyle\sum_v w_v-\frac{\sum_v w_v^2}{\sum_v w_v}\Bigr)
\quad\Longrightarrow\quad
\hat\tau_d^2=\frac{\bigl[Q-(m-1)\sigma^2\bigr]_+}{\sum_v w_v} ,
\end{equation}
(the implementation drops the $\sum w^2/\sum w$ term, a conservative simplification).
Two selection-aware modifications are essential:

\begin{enumerate}[leftmargin=2em]
\item \emph{Slab, not pooled.} Under a sparse (spike-and-slab) truth, estimating
$\tau_d^2$ over \emph{all} candidates dilutes the slab with nulls and overshrinks every
admitted coefficient. Consistency with the two-groups gate requires estimating
$\tau_d^2$ over \emph{admitted} vertices only: the gate decides existence, the prior
describes the nonnull component. Mild upward selection bias in $\hat\tau^2$ translates
into mild undershrinkage of a set whose false fraction is already bounded by $q$.
\item \emph{A floor, and reversion off support.} An information-starved vertex with
near-zero $x^\top Wx$ takes an effectively unregularized height; its error propagates
into neighbors' residuals, inflates nearby candidate $z$'s, opens the gate into the
damage, and the loop diverges (observed: heights of order $10^{11}$). A shrinkage floor
(equivalently a height sanity clamp resetting to $\mu_{\mathrm{lin}}$) removes the
feedback. Separately, any active vertex whose star contains no data is reverted to the
prior mean $0$: off the data support the deliverable is exactly the internal model.
The prior mean $0$ is the correct one for a \emph{correction} surface; it is what makes
the empty-space behavior auditable.
\end{enumerate}


\section{The admission null: conditional, composite, and permutation formulations}\label{sec:nulls}

The gate of Section~\ref{sec:gate} tests each candidate surplus against a
null distribution, and the choice of that null is a modeling decision with
measurable consequences. Three formulations are now implemented and
characterized. Throughout, the statistic template is
$z_v = (\hat\beta_v - \hat\mu_v)/\mathrm{sd}_v$ feeding the empirical-null
and prefix machinery of Section~\ref{sec:lfdr} unchanged; the formulations
differ only in $\mathrm{sd}_v$ and in its calibration.

\subsection{Conditional null (baseline)}
$\mathrm{sd}_v^2 = (1/X_v + \tfrac14 s_p)\hat\varphi$ with $X_v$ the
star information and $\hat\varphi$ the global Pearson dispersion. This is
the statement that, given the full latent state, $p$ defines pure noise.
The statement is true, and it is the wrong conditioning for the decision:
the vertex never observes the pool state, so under this null pool-level
overdispersion is read as location signal, and structure is admitted at
the wrong address. Under this formulation the empirical-null clamps
($\hat\sigma_0\in[1,1.4]$, $\hat\pi_0\ge0.5$, the depthwise $\tau^2$
floor) are load bearing: removing the $\tau$ floor triggers a shrink,
residual-signal, re-admission spiral (measured on the synthetic
benchmark: mesh $2127\to4697$, held-out $20.7\to27.8$).

\subsection{Composite null (analytic pool-variance form)}
The declared model has two signal layers, surface and pool, and a surface
hypothesis must be tested against the composite null that includes the
pool layer:
\[
\mathrm{sd}_v^2 = \Bigl(\frac{1}{X_v}
 + \sigma^2_{\mathrm{gate}}\,\frac{X^{(2)}_v}{X_v^2}
 + \tfrac14 s_p\Bigr)\hat\varphi,
\qquad X^{(2)}_v = \sum_i \tilde w_i^2 \varphi_v(z_i)^2,
\]
with $\sigma^2_{\mathrm{gate}}$ measured by the per-size instrument
(Section~\ref{sec:physics}). Fit weights are untouched; only admission
requires cross-pool replication. The identification is explicit: pool
noise averages at rate $K_{\mathrm{eff}}^{-1/2}$ with
$K_{\mathrm{eff}} = X_v^2/X^{(2)}_v$ the effective pool count, while a
shared surface shift survives averaging; a star dominated by one pool has
divergent null variance, the continuous form of a minimum-effective-pools
rule. Unreplicated variation is not discarded; it is routed to the pool
layer, which is its correct address. Under this null the clamps were
measured to be compensation for the conditional formulation: the $\pi_0$
floor is inert to the printed digit, the $\sigma_0$ box is near-inert, and
the $\tau$ floor is removable and mildly beneficial.

\subsection{Permutation null (exact conditional-on-effects form)}
The composite null has an assumption-free exact form. Take the null to be:
no spatial signal within the star; residual variation is pool specific.
The pool effects are then nuisance parameters to condition on, not draws
from a law to be modeled: the null distribution of the candidate statistic
is obtained by permuting the observed (residual, variance) pairs across
the star's hat-weight slots and recomputing the statistic. The null then
carries the realized tails, lumps, size structure, and issuer effects with
no mixing-law model at all, and any computation shared by the real and
permuted scores (baselines, fit-state transients) cancels identically.
Measured: on the noise-physics benchmark of Section~\ref{sec:certnull}
this matches the analytic composite gate (6 versus 7 certified-false deep
faces against the baseline's 997) with zero distributional assumptions and
slightly better held-out deviance (16.25 versus 16.8); with weak clean
Gaussian noise it exposes fit-state spatial transients as signal (66 on
the v3 benchmark); on real data it is over-conservative (730 vertices)
because all within-star structure, including genuine sub-face surface, is
priced into the null. It is the correct certification instrument under
realistic noise, and its real-data conservatism is the self-masking cost
of the composite stance made explicit. (Default-gate counts in the table
below are the post-fix values 969/996; an earlier draft carried 886/997
from before the dispersion-denominator correction of
Section~\ref{sec:onelaw}. The composite counts 2/7 are unchanged by that
fix.)

\subsection{The two configurations, precisely}
\begin{center}
\begin{tabular}{p{3.4cm}p{4.6cm}p{5.4cm}}
\toprule
component & conditional (baseline) & composite (principled) \\
\midrule
extra variance is & signal at the surface address & noise for the surface
test; signal for the pool layer \\
clamps & load bearing & measured inert or removable \\
single-pool stars & admissible & not certifiable \\
certified-null v3 deep faces & 969, false with certainty & 2 \\
noise-physics v4 deep faces & 996 & 7 analytic; 6 permutation \\
real held-out deviance/pool & 24.54 & 10.00 \\
marginal NLL per pool & 3.351 & 3.509 \\
in-zone info-weighted deviance & 33.68 & 11.12 \\
\bottomrule
\end{tabular}
\end{center}

The three evaluation metrics disagree by construction
(Section~\ref{sec:metrics}): on a single snapshot the fine-scale structure
is a bundle of genuine sub-face surface and pool-specific variation, and
the two configurations are the two set-identification bounds, the
conditional spending the bundle and the composite pricing it as null. The
measured $(\sigma_{\mathrm{gate}}, q)$ frontier shows the $q$ dial
dominating the $\sigma$ dial: keeping the null fully priced while
admitting more liberally reaches marginal NLL $3.442$ at $51$
certified-false faces, against $3.351$ at $886$ and $3.509$ at $2$.

\section{Residual variation and pool attribution}\label{sec:attribution}

After the gate stops, let $\hat p_i=\sigma(\eta_{0,i}+\hat g(x_i))$,
$\hat w_i=n_i\hat p_i\hat q_i$, and define Pearson residuals
$r_i=(k_i-n_i\hat p_i)/\sqrt{\hat w_i}$. Write the total unmodeled log-odds deviation
as $D_i=\{g(x_i)-\hat g(x_i)\}+u_i$, and decompose its variance as
\[
\Var(D)=\underbrace{\sigma_{\mathrm{sh}}^2}_{\text{spatially coherent miss}}
+\underbrace{\sigma_{\mathrm{id}}^2}_{\text{idiosyncratic }(u,\ \text{frailty, path effects})}.
\]

\begin{figure}[ht]
\centering
\includegraphics[width=\textwidth]{fig18_residual_attribution.png}
\caption{Conceptual map of the post-fit attribution layer. Binomial variation is known from exposure and fitted probability; coherent residual structure is potentially recoverable by a richer spatial fit; repeated observations identify persistent and white pool components through their different temporal behavior.}
\label{fig:attribution-concept}
\end{figure}

\subsection{Total moment}
Start from the exact binomial decomposition, valid for any $D_i$: with true
$p_i=\sigma(\logit\hat p_i+D_i)$,
\[
\E(k_i-n_i\hat p_i)^2=n_ip_iq_i+n_i^2(p_i-\hat p_i)^2 .
\]
Expanding both terms at $\hat p_i$ (delta method; $dp/d\eta=pq$,
$d(pq)/d\eta=(1-2p)pq$) yields the per-observation identity that drives all the
estimators of this section.

\begin{proposition}[delta-method identity, binomial]\label{prop:total}
\begin{equation}\label{eq:deltaid}
\E[r_i^2]-1
\;=\;
\underbrace{\hat w_i\,D_i^2}_{\text{energy}}
\;+\;
\underbrace{(1-2\hat p_i)\,D_i}_{\text{leakage}}
\;+\;O\!\bigl(D_i^2\bigr)\ \text{terms not scaled by }\hat w_i .
\end{equation}
Consequently, if the leakage term is negligible (see below), with $\mathrm{edf}$ the
effective degrees of freedom of the fit,
\begin{equation}\label{eq:totalmoment}
\widehat{\Var(D)}_{\mathrm{tot}}
=\frac{\sum_i r_i^2-(N-\mathrm{edf})}{\sum_i \hat w_i}
\end{equation}
estimates the $\hat w$-weighted energy $\E_{\hat w}[D^2]$. The level is estimated
not from \eqref{eq:deltaid} but at first-moment efficiency,
\begin{equation}\label{eq:level}
\hat L=\frac{\sum_i\sqrt{\hat w_i}\,r_i}{\sum_i\hat w_i},
\qquad \mathrm{se}(\hat L)\approx\Bigl(\textstyle\sum_i\hat w_i\Bigr)^{-1/2},
\end{equation}
and the energy moment is corrected by
$\hat L\cdot\sum_i(1-2\hat p_i)/\sum_i\hat w_i$ when $\hat L$ is materially nonzero.
\end{proposition}

\emph{Estimand caveat.} The identity identifies a $\hat w$-weighted energy.
The gate refines where information lives, so the residual miss $D^2$
concentrates where $\hat w$ is small; $\hat w$ and $D^2$ are negatively
correlated by construction and the weighted energy understates the unweighted
energy of the miss. Feeding the resulting $\hat\sigma^2$ into the DL weights
as a homogeneous variance is a modeling substitution the fixed-point argument
does not license; it is adopted for its variance-control behavior, verified
empirically. The neglected $O(D^2)$ terms not scaled by $\hat w$ are also
relatively largest for thin pools near the information floor; again where
the miss lives.

\begin{remark}[score annihilation of the leakage]\label{rem:leakage}
At an exact ML fixed point the leakage vanishes: the score equations state
$\sum_i(k_i-n_i\hat p_i)\varphi_v(x_i)=0$ for every active vertex, and since the hat
functions form a partition of unity, $\sum_v\varphi_v\equiv 1$, summing gives
$\sum_i(k_i-n_i\hat p_i)=0$, i.e.\ $\sum_i\hat w_iD_i\approx 0$ to first order.
Regularization breaks this exactly: EB shrinkage attenuates coefficients, the gate
refuses sub-threshold structure, and the outer IRLS terminates finitely, so a
systematic mean residual survives and contaminates \eqref{eq:totalmoment} when $D$ is
large. The level statistic \eqref{eq:level} is the right instrument for it: estimating a
first moment from $r^2$ (e.g.\ by regressing $r^2-1$ on $(\hat w,1-2\hat p)$) costs
orders of magnitude in efficiency and was found empirically to return pure noise even
with $1-2\hat p$ spanning $\pm0.98$, while \eqref{eq:level} resolved the same
quantity to $\pm 0.001$ log-odds. (That precision assumes independent residuals; within-pool frailty correlation implies a design effect that inflates it, so the quoted se is a lower bound.) In sample, the weighted score of exact ML pins
$\hat L$ at zero, so $\hat L$ measures precisely the level that regularization leaves
uncorrected (a direct gauge of EB-plus-gate absorption; it should shrink as the
fixed-point iteration of Section~\ref{sec:posthoc} converges, and it grew
monotonically as the gate was tightened in the $q$-sweep, as designed). Two regime
notes: with $p$ spread widely, $\sum_i(1-2\hat p_i)\approx 0$ and the leakage
correction self-cancels; at hazard scale ($p\sim1\%$, $1-2\hat p\approx1$) it does
not, and a level of $-0.004$ shifts a post-gate energy of $0.01$ by roughly $20\%$,
which is where the correction earns its keep.
\end{remark}

\begin{remark}[validity boundary]
The first neglected contribution to \eqref{eq:deltaid} is a skewness term of order
$\hat w_i(1-2\hat p_i)\,\E D^3$: with a large negatively skewed miss the moment
under-reads (observed empirically on a deliberately degraded fit). No distributional
assumption on $D$ is made anywhere above, the output is directly in log-odds units
(the units the DL weights and the Gamma prior consume), and the small-to-moderate
miss regime after the gate is exactly where the expansion is sharp.
\end{remark}

The following exact alternative covers the large-$D$ regime at the price of a
distributional assumption and a rate-scale detour; it is retained as a cross-check
rather than as the primary statement.

\begin{proposition}[exact Poisson--lognormal inversion]\label{prop:lognormal}
Suppose $p$ is small, $k\mid D\sim\mathrm{Poisson}(\lambda e^{D})$ with
$D\sim N(0,s)$, and the fit is mean-matched, $\hat\lambda=\lambda\E e^{D}=\lambda
e^{s/2}$. Then, with $r=(k-\hat\lambda)/\sqrt{\hat\lambda}$ and information
$\hat w=\hat\lambda$,
\begin{equation}\label{eq:exact}
m \;\equiv\; \frac{\E[r^2]-1}{\hat w}
\;=\; \frac{\E (k-\hat\lambda)^2-\hat\lambda}{\hat\lambda^{2}}
\;=\; e^{s}-1
\qquad\Longleftrightarrow\qquad
s=\log(1+m).
\end{equation}
\end{proposition}

\emph{Register.} The inversion is exact \emph{under its mean-matching premise}
$\hat\lambda=\lambda e^{s/2}$, which the pipeline's penalized, gated fit does
not satisfy (that is the content of the broken-score-annihilation remark), so
this proposition serves as a cross-check under an idealization rather than a
property of the produced fit.
\begin{proof}[Calculation]
By the conditional variance decomposition,
$\E(k-\hat\lambda)^2=\E[\Var(k\mid D)]+\Var(\E[k\mid D])
=\E[\lambda e^{D}]+\Var(\lambda e^{D})
=\lambda e^{s/2}+\lambda^2\bigl(e^{2s}-e^{s}\bigr)$,
using $\E e^{D}=e^{s/2}$ and $\Var(e^{D})=\E e^{2D}-(\E e^{D})^2=e^{2s}-e^{s}$.
Subtracting $\hat\lambda=\lambda e^{s/2}$ leaves
$\lambda^2(e^{2s}-e^{s})=\hat\lambda^{2}(e^{s}-1)$ since
$\hat\lambda^{2}=\lambda^{2}e^{s}$; dividing by $\hat\lambda^{2}$ gives
\eqref{eq:exact}. The empirical statistic
$(\sum_i r_i^2-(N-\mathrm{edf}))/\sum_i\hat w_i$ is the sample version of $m$, so the
missed variance is recovered as $\hat s=\log(1+\hat m)$.
\end{proof}
Expanding $\log(1+m)=m-\tfrac{m^2}{2}+\dots$ shows Proposition~\ref{prop:total} is the
first-order term; at the development magnitudes ($s\approx0.12$, i.e.\ missed sd
$\approx0.35$ log-odds) the correction is $\sim -8\%$ of $m$ and grows with $s$. In the
intended deployment the gate runs first, $s$ is small, and the delta-method
regression of Proposition~\ref{prop:total} is the recommended estimator; the exact
inversion is a cross-check where its assumptions plausibly hold.

\begin{remark}[post-selection bias and the holdout moment]\label{rem:holdout}
Formula \eqref{eq:totalmoment} takes each fitted degree of freedom to absorb one unit
of noise. Gate-admitted vertices are selected \emph{for} large coefficients: under the
null, $\E[z^2\mid |z|>t]=1+t\phi(t)/\{1-\Phi(t)\}\approx 4$--$9$ for $t\in[2,3]$, so an
adaptively selected fit absorbs several noise units per admitted vertex and
\eqref{eq:totalmoment} is biased \emph{low}, by an amount matching the observed
shortfall ($\approx \mathrm{edf}\times 3/\sum\hat w$) in development. The clean repair
is a holdout moment: evaluate $r$ on data the fit never saw (a temporal holdout in
production); then $\mathrm{edf}=0$, no selection adjustment exists, and
\eqref{eq:exact} applies directly.
\end{remark}


\begin{figure}[h]
\centering
\includegraphics[width=\textwidth]{fig3_residuals.png}
\caption{The estimator's core identity, visualized. Left: Pearson residual distributions with ($\sigma_u=0.30$) and without frailty. Right: binned $\mathbb{E}[r^2]-1$ against per-observation information $w$: with frailty the excess is linear in $w$ with slope $\approx\sigma_u^2=0.09$ (dashed reference); without frailty it is flat near zero, with the mild positive drift at large $w$ reflecting residual surface miss under the density cap. Single seed, $4\times$ data.}
\label{fig:excess}
\end{figure}

\subsection{Shared and idiosyncratic components}
Partition observations into small cells $C$ (side just above the finest mesh scale) and
let $D_i=\Delta_{C(i)}+\varepsilon_i$ with $\Delta_C$ the cell-coherent miss and
$\varepsilon_i$ idiosyncratic across units.

\begin{proposition}[pairwise moment]\label{prop:pair}
For $i\neq j$ in the same cell and different units,
$\E[r_ir_j]=\sqrt{\hat w_i\hat w_j}\,\sigma_{\mathrm{sh}}^2+o(\cdot)$, and
\begin{equation}\label{eq:pair}
\hat\sigma^2_{\mathrm{sh}}
=\frac{\sum_C\bigl[(\sum_{i\in C}r_i)^2-\sum_{i\in C}r_i^2\bigr]
-\text{(same-unit pair terms)}}
{\sum_C\bigl[(\sum_{i\in C}\sqrt{\hat w_i})^2-\sum_{i\in C}\hat w_i\bigr]
-\text{(same-unit pair terms)}}
\end{equation}
is unbiased for the cell-scale-coherent miss; the same construction over within-unit
(same pool, different month) pairs estimates the unit-persistent component
$\sigma_{\mathrm{pool}}^2$ (frailty plus persistent path effects), and
$\hat\sigma^2_{\mathrm{id}}=\widehat{\Var(D)}_{\mathrm{tot}}-\hat\sigma^2_{\mathrm{sh}}$.
\end{proposition}
\begin{proof}[Calculation]
$r_i\approx\sqrt{\hat w_i}\,D_i+\nu_i$ with $\nu$ the standardized binomial noise,
independent across $i$; cross terms vanish and
$\E r_ir_j=\sqrt{\hat w_i\hat w_j}\,\E[\Delta_C^2]$ for distinct units. Summing
$(\sum r)^2-\sum r^2=\sum_{i\neq j}r_ir_j$ over cells and normalizing by the analogous
weight sums gives \eqref{eq:pair}.
\end{proof}

\begin{remark}[the two estimands]\label{rem:estimands}
The pairwise family (and the equivalent variogram nugget-difference estimator) is blind
to miss that is incoherent at cell scale. After a \emph{successful} fit, the remaining
error is P1 interpolation truncation, which oscillates within faces and averages out at
cell scale; the pairwise estimator then correctly reads $\approx 0$ while the total
moment reads the truncation energy. The two are therefore not competitors for one
estimand: the total moment measures everything missed; the pairwise moment measures the
\emph{findable} part (recoverable by further refinement or data); their difference is
resolution-limited truncation plus idiosyncratic variance. In the Monte Carlo
comparison on cosine data with cohort geometry, the total moment had bias $-0.011$,
RMSE $0.038$; cross-validation was unbiased at triple the RMSE; the pairwise and
nugget estimators tracked truth when the fit had failed and read near zero when it had
succeeded, exactly as the estimand analysis predicts.
\end{remark}

\subsection{The report}
The deliverable diagnostic after the gate stops is:
$\hat\sigma_u^2$ from the leverage-corrected pool moment (REML-like in the balanced limit; Paule--Mandel-type moments and REML diverge under weight heterogeneity, which is precisely our regime)
(Section~\ref{sec:doubledip}); the total energy from the $(\hat w,\,1-2\hat p)$
regression of Proposition~\ref{prop:total}, post-hoc centered with its trace or
bootstrap correction; the findable component as total minus idiosyncratic,
cross-checked by the pairwise moment. A holdout evaluation remains available as
insurance, though at production sample sizes the selection absorption it guards
against was measured to be negligible and is in any case covered by the bootstrap
correction.
Interpretation: findable $\Rightarrow$ refine or add data; idiosyncratic
$\Rightarrow$ the missing structure is not in these coordinates (burnout state, pool
frailty), pursue it with an additional state variable or unit-level expression, not
with more triangles. Paired with Proposition~\ref{prop:fdr}, the procedure carries a
two-sided internal accounting, with the FDR side read under the conditional-on-mesh scope of the remark following Proposition~\ref{prop:fdr} and the conditioning table of Section~\ref{sec:conditionality}: at most $q$ of the admitted structure is false in the per-round Bayes sense, and the
undiscovered remainder is measured and attributed. The idiosyncratic component is not
discarded: Section~\ref{sec:posthoc} converts it, pool by pool, into shrunken
posterior attributions.


\subsection{Credibility-weighted pool attribution and the variance fixed point}\label{sec:posthoc}

After the surface is fitted and the variance components are measured, each pool's
realized error admits a posterior attribution into systematic and idiosyncratic parts,
and the attribution changes the estimate for that pool going forward. This section
makes the update exact and identifies the fixed point it satisfies.

\subsection{Credibility form of the posterior}
On the working scale, model observation $i$ of pool $P$ as
$y_i=g(x_i)+u_P+\varepsilon_i$ with $u_P\sim N(0,\sigma_u^2)$ independent of
$\varepsilon_i\sim N(0,1/w_i)$. Given the fitted surface, the residuals
$e_i=y_i-\hat g(x_i)$ over the pool's observed months carry information
$W_P=\sum_{i\in P}w_i$, and the posterior of $u_P$ is Gaussian with
\begin{equation}\label{eq:blup}
\hat u_P=\E[u_P\mid e]
=\frac{W_P}{W_P+1/\sigma_u^2}\;\bar e_P^{\,w},
\qquad
\Var(u_P\mid e)=\frac{1}{W_P+1/\sigma_u^2},
\end{equation}
where $\bar e_P^{\,w}=\sum_{i\in P}w_ie_i/W_P$. This is the BLUP of mixed-model theory
and the B\"uhlmann credibility premium \cite{buhlmann1967,henderson1975}: the observed
pool error is split between idiosyncratic signal and noise in the ratio
$W_P\sigma_u^2:(W_P\sigma_u^2+1)$, so a small pool with two lucky months is shrunk
essentially to zero while a large, persistently deviating pool keeps most of its
residual. The same formula one level up, applied to cross-pool cell means with
$\sigma_{\mathrm{sh}}^2$ in place of $\sigma_u^2$, is the post-hoc revision of the
surface itself; the moment estimators of Section~\ref{sec:attribution} supply exactly the
variance components these weights require.


\begin{figure}[h]
\centering
\includegraphics[width=\textwidth]{fig4_blup.png}
\caption{Pool attribution on simulated data ($\sigma_u=0.30$, $4\times$). Left: credibility posterior $\hat u_P$ against true frailty, colored by pool information; corr $0.78$. Right: empirical shrinkage (regression of $\hat u_P$ on $u_P$ within information bins): low-information pools recover $\sim$40--65\% of their frailty and well-measured pools $\sim$95\%, so a thin pool cannot top a ranking on two lucky months. Single seed.}
\label{fig:blup}
\end{figure}

\subsection{The fixed point}
The shrinkage weights in \eqref{eq:blup} require $(\sigma_{\mathrm{sh}}^2,\sigma_u^2)$;
the moment estimators of these components are computed from residuals against the
fitted effects; the fitted effects depend on the weights. The joint solution
$(\hat g,\hat u,\hat\sigma_{\mathrm{sh}}^2,\hat\sigma_u^2)$ is therefore a fixed point,
and it is the familiar one: the stationary point of the (restricted) likelihood of the
mixed model. Alternating (i) posterior effects given variances with (ii) variance
re-estimation given effects resembles an EM alternation, but the moment update is not an M-step and no monotone-ascent guarantee is claimed; with the surface refit and the gate's selection inside the loop, existence and stability of the fixed point are empirical observations in this document,
so the fixed point exists and the iteration is stable under standard conditions.
The pipeline of Sections~\ref{sec:glm}--\ref{sec:overdisp} is a truncated version:
one-step scoring, DL moment \eqref{eq:dl}, DL reweighting \eqref{eq:dlw}, refit is two
Picard sweeps of the same map, and the moment-iterated variant (re-solving the DL
identity at the updated weights until convergence) is the Paule--Mandel estimator
\cite{paule1982}. Running the alternation to convergence is the natural finishing move
of the algorithm and is how the implementation should terminate when the pool layer is
in use.

\subsection{Exact conjugate form on the hazard scale}
The Gaussian working-scale posterior is an approximation; on the hazard scale the
update is exact. Let the pool carry a multiplicative frailty $\theta_P$ on the hazard,
$\theta_P\sim\mathrm{Gamma}(\alpha,\beta)$ with $\E\theta=\alpha/\beta=1$, and let the
month-$t$ events be (small-$p$) Poisson with mean $\theta_P\,\Lambda_{P,t}$, where
$\Lambda_{P,t}=n_{P,t}\hat p_{P,t}$ is the fitted exposure. Then
\begin{equation}\label{eq:gammapost}
\theta_P\ \Big|\ \{k_{P,t}\}\ \sim\
\mathrm{Gamma}\Bigl(\alpha+\textstyle\sum_t k_{P,t},\ \ \beta+\sum_t\Lambda_{P,t}\Bigr),
\qquad
\E[\theta_P\mid\cdot]=\frac{\alpha+\sum_t k_{P,t}}{\beta+\sum_t\Lambda_{P,t}} ,
\end{equation}
closed form in cumulative counts and cumulative fitted exposure, with no logit
approximation; matching $\Var\theta=1/\alpha$ to the measured heterogeneity gives
$\alpha=1/(e^{\sigma_u^2}-1)$ under the lognormal reading of
Proposition~\ref{prop:lognormal}, or $\alpha=1/\sigma_u^2$ to second order. Two
consequences. First, the sufficient statistic $\sum_t\Lambda_{P,t}$ is cumulative
fitted exposure, i.e.\ precisely the burnout state of the population-dynamics
description (Section~\ref{sec:dynamics}): the posterior update of a pool's frailty
\emph{is} the selection functional, one state variable doing double duty for
composition drift and for post-hoc attribution. Second, \eqref{eq:gammapost} exhibits
the credibility structure exactly: the posterior mean is a convex combination of the
prior mean $1$ and the pool's empirical hazard ratio, with weight
$\sum\Lambda/(\beta+\sum\Lambda)$.

\subsection{Dynamics, and what the update is epistemically}
If the pool effect decays, $u_{P,t+1}=\phi\,u_{P,t}+\text{innovation}$, the static
posterior \eqref{eq:blup} is the steady state of a per-pool Kalman filter whose gain is
again a credibility weight; the optimal $h$-step forecast of a pool's error is the
current posterior mean scaled by $\phi^h$, which is the reliability-times-persistence
form measured directly by the frame-attachment panel of Section~\ref{sec:dynamics}.
Finally, the epistemic status of $\hat u_P$ differs from that of the surface: admitted
surface structure is FDR-certified discovery; $\hat u_P$ is a shrunken prediction with
posterior variance \eqref{eq:blup}, licensed by compound estimation (the
hyperparameters are pooled over thousands of pools, so no per-pool selection is
involved) rather than by a test. Ranked by $\hat u_P$, the pool layer is a selection
engine for unit-level expression (specified pools, IO/PO) in which the shrinkage
guarantees that position size in a pool's ``story'' is automatically limited by the
information supporting it.

\subsection{Consistency of the attribution boundary}\label{sec:leverage}

Scaling experiments (sample size quadrupled; FDR target swept over
$q\in[0.02,0.5]$) show a clean split. Estimators whose truth is defined relative to
the fitted surface are consistent: the energy tracks its (shrinking) truth within a
few percent at every scale and every $q$, the decomposition identity
$\text{energy}=\text{post-hoc}+\text{explained}$ holds throughout, and the pool
posterior calibration stays at $1.0$. The $q$-sweep additionally traces the expected
detection U-curve: true miss falls from $q=0.02$ to a minimum near $q\in[0.1,0.25]$
and rises again at $q=0.5$ as the gate admits noise.

The internal \emph{boundary} of the decomposition, however, carries an $O(1)$ bias
that does not shrink with $N$ and is invariant to $q$: $\hat\sigma_u^2$ reads
$\approx 9\%$ low in sample at every scale. Attribution experiments exonerate the two
natural suspects. An oracle test (pairwise moment at the \emph{true} systematic $p$,
no fitting) is unbiased to $\pm1\%$ with and without known offsets, so the
delta-method expansion is not the cause; and removing the known per-pool offset from
the generative model leaves the bias unchanged. The cause is fit
\emph{leverage}: pools-per-vertex-star stays $O(1)$ as the gate refines in proportion
to $N$, so the surface permanently absorbs a fraction of each pool's $u$. In sample
the moment is deflated by roughly $(1-h_P)^2$ in the pool's leverage $h_P$; on
held-out pools the same mechanism inverts, the surface missing exactly that pool's
contribution along its own track, inflating the moment by roughly $(1-h_P)^{-2}$
(measured: $0.083$ and $0.126$ around a truth of $0.090$). That the holdout
contamination is pool-own leverage and not spatial is proved by the cross-pool
within-cell moment on the same holdout residuals reading $\approx 0$. The geometric
mean of the bracket ($0.102$) is a crude leverage-robust combination, exact if $h_P$
were homogeneous; it is not, being concentrated in large pools.

\subsubsection*{Analytic leverage correction}
On the working scale after the final expansion, $y=\Phi h+Zu+\varepsilon$ with
$\Phi$ the $N\times V$ hat design, $Z$ the pool indicators,
$u\sim(0,\sigma_u^2 I_P)$, $\varepsilon\sim(0,V)$, $V=\mathrm{diag}(1/w_i)$ the
\emph{binomial} noise variances. The surface is the linear smoother $S=\Phi
A_\Lambda^{-1}\Phi^\top \widetilde W$ with $A_\Lambda=\Phi^\top\widetilde W\Phi+
\Lambda$ and $\widetilde W$ the (DL) fitting weights; residuals $e=(I-S)y$, and if
the admitted surface spans $g$,
\[
e=(I-S)Zu+(I-S)\varepsilon,\qquad M:=I-S .
\]
The pairwise statistic is $T=\sum_P\bigl[(\sum_{i\in P}r_i)^2-\sum_{i\in P}r_i^2\bigr]$
with $r_i=s_ie_i$, $s_i=\sqrt{w_i}$, i.e.\ a sum over same-pool pairs
$\sum_P\sum_{i\neq j\in P}s_is_j\,e_ie_j$. Taking expectations termwise,
\begin{equation}\label{eq:master}
\E[e_ie_j]=\sigma_u^2\,(Mz_{P})_i(Mz_{P})_j+(MVM^\top)_{ij}
\qquad (i,j\in P,\ i\neq j),
\end{equation}
up to cross-pool leak-in terms $(Sz_Q)$, second order and omitted. Hence the exactly
unbiased estimator is
\begin{equation}\label{eq:levcorr}
\hat\sigma^2_{\mathrm{corr}}
=\frac{T-C_{\mathrm{noise}}}{D_{\mathrm{sig}}},\qquad
D_{\mathrm{sig}}=\sum_P\bigl(a_P^2-b_P\bigr),\quad
C_{\mathrm{noise}}=\sum_P\sum_{i\neq j\in P}s_is_j\,(MVM^\top)_{ij},
\end{equation}
where, with the pool-leverage profile
$h_i^{(P)}=\varphi_i^\top A_\Lambda^{-1}c_P$, $c_P=\Phi^\top\widetilde W z_P$,
\[
a_P=\sum_{i\in P}s_i\bigl(1-h_i^{(P)}\bigr),\qquad
b_P=\sum_{i\in P}s_i^2\bigl(1-h_i^{(P)}\bigr)^2 .
\]
The naive estimator replaces $D_{\mathrm{sig}}$ by $\sum_P[(\sum s_i)^2-\sum s_i^2]$
and $C_{\mathrm{noise}}$ by zero, and both replacements bias it \emph{down} through
two distinct channels: the smoother reproduces part of $z_P$, shrinking
$1-h_i^{(P)}$ (the signal channel), and same-pool observations are co-located, so
$S_{ij}>0$ between them and $(I-S)\varepsilon$ acquires negative within-pool
covariance (a noise channel the pure-leverage heuristic misses; this is why the
geometric mean of the in-sample/holdout bracket under-corrects). All ingredients are
quadratic forms in $A_\Lambda^{-1}$: with $\widetilde W=V^{-1}$, $\Lambda=0$ one has
$(MVM^\top)_{ij}=-\varphi_i^\top A^{-1}\varphi_j$ for $i\neq j$; in general
$(MVM^\top)_{ij}=-\varphi_i^\top A_\Lambda^{-1}\varphi_j(\tilde w_jv_j+\tilde
w_iv_i)+\varphi_i^\top A_\Lambda^{-1}BA_\Lambda^{-1}\varphi_j$ with
$B=\Phi^\top\widetilde WV\widetilde W\Phi$. Since each $\varphi_i$ has three
nonzeros, the cost is one factorization of $A_\Lambda$ plus local back-substitutions,
cheaper than a height sweep. To leading order, if $h_i^{(P)}\approx h_P$ over a
pool's observations, $D_{\mathrm{sig}}\approx(1-h_P)^2$ times the naive denominator,
recovering the deflation quoted above. Numerically, with $\Lambda=0$ the
correction over-shoots ($+17\%$ at $\sigma_u=0.1$, where the mesh is richest and
shrinkage strongest, since the unpenalized $S$ overstates absorption); with the EB
surplus ridge included, $\Lambda_{\mathrm{eff}}=C^\top\mathrm{diag}(1/\tau_d^2)C$
where $C$ maps heights to surpluses ($+1$ at $v$, $-\tfrac12$ at each parent), the
corrected estimator lands within $\pm4\%$ across $\sigma_u\in[0.1,0.5]$
($+4\%$, $+1.2\%$, $-4\%$ at $\sigma_u=0.1,0.3,0.5$, against naive biases of
$-18\%$ to $-10\%$), the residual at the top of the range being the delta-method
and cross-pool leak-in leftovers. Substituting
$\hat\sigma^2_{\mathrm{corr}}$ back into the DL weights and the $\kappa$'s makes the
fixed point a REML-type score equation of the \emph{joint} surface-plus-pools mixed
model, with the surface's absorption in the accounting. Validity: the identity lives
on the working scale (delta-method domain, $\pm2\%$ to $\sigma_u=0.7$ by the oracle
sweep), treats $S$ as fixed (the gate's selection component remains the bootstrap's),
and does not see the logistic-saturation effect that breaks the holdout bracket at
$\sigma_u\gtrsim0.5$.

The certified consistent estimator is bootstrap calibration of the \emph{full}
pipeline: simulate at the fitted model including the surface refits, measure the map
$\sigma^2\mapsto\E[\hat\sigma^2]$, and invert it at the data estimate; this
reproduces the leverage heterogeneity exactly and is an offline job ($B$ complete
fits). Two structural mitigations reduce the need for it: capping any single pool's
weight share within a vertex star bounds $h_P$ directly, and longer observation
histories shrink leverage mechanically (production residence times of $60$-plus
months against the simulated $4.5$). The summary statement for governance: the
decomposition's \emph{totals} are consistent; its internal boundary is biased by fit
leverage, bracketed by the in-sample/holdout pair, and certified by pipeline
bootstrap calibration.

\subsection{Post-hoc centering and the double-dip correction}\label{sec:doubledip}

The moments of Section~\ref{sec:attribution} sharpen when centered at the post-hoc
probability $\tilde p_i=\sigma(\eta_{0,i}+\hat g(x_i)+\hat u_{P(i)})$: the
pool-persistent component leaves $D$, the skewness boundary of
Proposition~\ref{prop:total} recedes, and the resulting statistic is the innovation
variance the dynamic pool filter consumes. The cost is double dipping: $\hat u_P$ was
fit to the same counts entering the residuals, and a posterior mean absorbs noise, so
post-hoc Pearson residuals are underdispersed by construction. The deficit, however,
is estimable and correctable rather than avoidable-only.

\emph{Linear layer: trace correction $=$ REML.} The credibility posterior is a
rank-one linear smoother per pool, so under a correct model
\[
\E\Bigl[\sum_i \tilde r_i^{\,2}\Bigr]=N-\sum_P\kappa_P,
\qquad
\kappa_P=\frac{W_P}{W_P+1/\sigma_u^2},
\]
i.e.\ the effective degrees of freedom of the pool layer is the sum of its own
shrinkage factors. Correcting the moment by $\sum_P\kappa_P$, and noting that
$\kappa_P$ depends on $\sigma_u^2$ so that the corrected moment equation is a
one-dimensional monotone fixed point, reproduces the ML-to-REML adjustment
for the variance component in the balanced Gaussian limit: there, and only
there, the double-dip correction coincides with REML.

\emph{Nonlinear layers: bootstrap $=$ covariance penalty.} Where no trace formula
exists (the Gamma posterior is nonlinear in counts; the gate's selection absorption;
the stability clamps), simulate $B$ replicates from the fitted model, push each
through the identical estimator pipeline, and subtract the bootstrap-null mean of the
statistic. This is Efron's covariance-penalty estimate of generalized degrees of
freedom \cite{efron2004pred,ye1998} and corrects double dipping, selection, and
shrinkage uniformly. Agreement between trace and bootstrap on the linear layer is a
standing consistency check; material divergence on real data indicates some component
absorbing more than its linear-theory share.

\emph{Where the trace genuinely fails.} The linearized trace can miss absorption even
when the posterior-mean smoother is linear in the counts, because Pearson-type
statistics are also standardized by an estimated scale: in $(k-\hat\theta\Lambda)^2/
(\hat\theta\Lambda)$ the denominator co-moves with the numerator (a hot pool inflates
both), an additional negative coupling invisible to any mean-smoother trace. Under
Gamma centering with a correct model, the naive excess read $-0.0255$; the linearized
trace removed only $70\%$ of the deficit (residual $-0.0073$, implied df $1489$);
the bootstrap read $-0.0004$ with implied df $2049\pm13$, a $38\%$ df gap
attributable to the standardization coupling. The rule is therefore two lines: on the
Gaussian working-scale credibility layer, trace and bootstrap agree exactly and the
trace is free; wherever the variance function, the Gamma form, or selection touches
the statistic, the bootstrap is the correct answer and the trace-bootstrap gap
measures what the shortcut misses.

\emph{Numerical check} (isolated pool layer, $2{,}000$ pools, histories $3$--$30$
months, $\sigma_u^2=0.09$): the naive post-hoc excess reads $-0.018$ where a correct
model should read $0$; trace-corrected, $-0.00005$; bootstrap-corrected
($B=200$), $+0.00016$, with bootstrap-implied null degrees of freedom $1515$ against
the analytic trace $1498$. The uncorrected variance fixed point converges to
$\hat\sigma_u^2=0.115$ (absorbed noise counted as signal); the REML fixed point to
$0.092$.

\begin{remark}[sequential centering, retained for one purpose]
An alternative that avoids double dipping by construction is prequential centering,
$\hat u_{P,t}$ computed from months before $t$ only. As an \emph{estimator} it is
dominated by the corrected full-information posterior, since it discards data the
correction can account for. It retains one distinct role: one-step-ahead innovations
have conditional mean zero under a correct model, so the leakage coefficient of the
regression in Proposition~\ref{prop:total}, computed on sequentially centered
residuals, is a pure misspecification (innovation-whiteness) test rather than a
regularization diagnostic.
\end{remark}

\subsection{Pool state and its estimators}\label{sec:poolstate}

The pool-state model is
\[
u_{P,t}=a_P+\eta_{P,t},\qquad a_P\sim(0,\sigma_a^2),\qquad
\eta_{P,t}\sim(0,\sigma_w^2)\ \text{i.i.d.\ across months.}
\]
Under it the within-minus-shared lag profile of
Section~\ref{sec:attribution} is flat at $\sigma_a^2$ for every $\ell\ge1$,
so the pairwise moment is exactly unbiased for the persistent component with
no lag attenuation and the leverage correction of
Section~\ref{sec:leverage} applies unchanged. The white component is the structurally hard parameter, for an
identifiable reason: it lives \emph{only} on the diagonal, and the diagonal
is where everything unidentified lands. The persistent component has
off-diagonal structure to be matched on (the flat lag profile); white noise
by definition has none, so every diagonal moment reads
$\sigma_a^2+\sigma_w^2+m(z)^2$ plus selection error, and estimation is a
matter of subtracting the other three with instruments that do not touch
$\sigma_w^2$. The estimator assembles three such pieces. The total comes
from the Paule--Mandel iteration
\[
\sum_i \frac{r_i^2}{1+w_i\sigma}\;=\;N-\operatorname{tr}S(\sigma)
\]
($\operatorname{tr}S$ under working weights plus EB ridge, Hutchinson probes
through the sparse CG; monotone, bisected, under a second). The persistent
part is the leverage-corrected pairwise moment. The miss part is a zero-lag
instrument: co-located (same-face) pairs from \emph{different} pools share
the local miss $m(z_i)m(z_j)$ and nothing else, no frailty across pools and
no white noise across independent draws, so the same-face cross-pool moment
$Z$ estimates the information-weighted local miss second moment up to the
within-face coherence of the $P_1$ interpolation error ($\kappa=15/16$ for
the isotropic bubble). Then
$\hat\sigma_w^2 = \text{PM total} - \hat\sigma_{a,\mathrm{corr}}^2 - Z$.
The identifiability boundary is the zero-lag twin of the earlier ones: miss
varying below the face scale \emph{and} below the monthly trajectory step is
observationally equivalent to white pool noise, on the diagonal and along
trajectories alike, so $Z$ recovers the resolved miss and sub-face amplitude
variation remains inside $\hat\sigma_w^2$ by convention.

\emph{Verification} (uniform surface unless noted; three to five seeds per
cell; truth total $0.09$ split as $\sigma_a^2/\sigma_w^2$). Persistent
component: scale-consistent ($\hat\sigma_{a,\mathrm{corr}}^2 =
0.085/0.090/0.088/0.090$ across $1$--$16\times$ at permanent share $1$;
$0.048/0.045/0.044/0.044$ at share $\tfrac12$), gate-insensitive ($0.0447$
vs.\ $0.0453$ at $q=0.05/0.25$), and falsification-clean (pure white and
no-signal both read $-0.0001\pm.0012$: no manufactured persistence). Its
chirp inflation ($0.117/0.054$) is the attribution semantics of
Section~\ref{sec:attribution} operating as intended. White component
(miss-adjusted, three seeds per cell): $0.050/0.043/0.047$ across
$4/8/16\times$ at the half mix (truth $0.045$; the $16\times$ spread is
$\pm.0003$), $0.092/0.088$ under pure white noise (truth $0.09$),
$0.004/0.003$ in the pure-permanent case, and, decisively, $0.004\pm.004$
on the chirp, where the instrument routes the large multi-scale miss into
$Z=0.019$ rather than into $\hat\sigma_w^2$. The residual error decomposes as follows, dominant
cause first. (i) \emph{Selection gap in the PM trace}, now quantitatively
attributed: the trace treats the gate-selected basis as fixed, and the
double-dip analysis measured the selection surplus at $\approx38\%$ of the
nominal edf; with edf $\approx660$ and the measured PM sensitivity
$\partial\sigma/\partial\mathrm{edf} = 1.9\times10^{-5}$ (from
$\sum_i w_i/(1+w_i\sigma)^2 \approx 5.3\times10^4$ at the half mix), the
predicted bias is $+0.0047$ against the observed $+0.0052$; essentially the
entire systematic error. The definitive measurement is instrumented (split
randomness fixing data locations while resampling outcomes, per-observation
fit dumps) but requires a fixed-point linearization of the working response,
since the IRLS response contains the fit and naive covariance double-counts
it; that measurement, and the edf substitution it licenses, are the queued
repair. A first attempt at the split-randomness measurement was run and
diagnosed rather than trusted: the outcome-only replicate stream redraws
pool frailty $(a_P,u)$ alongside the binomial count, so replicates differ in
their \emph{true} $p_t$, not only in binomial sampling noise: the observed
spread of $\hat p$ across replicates (median $0.023$) exceeds the
pure-binomial expectation at typical pool size (${\approx}0.020$) only
mildly at the median but by an order of magnitude in the upper tail,
consistent with a subset of pools re-drawing frailty. Efron's formula then
divides by the wrong (too-small) variance and the resulting edf is not
usable. The correct instrument needs a three-way split, locations, frailty,
and terminal binomial draw, each independently controlled; the two-way split
here isolates the mechanism but not yet the number. (ii) \emph{$\kappa$ approximation}: $15/16$ is exact for the
isotropic interpolation bubble under uniform within-face density; NVB
anisotropy, lumpy density, and post-shrinkage (rather than interpolation)
miss profiles perturb it in either direction, bounded by the small size of
$Z$ itself. (iii) \emph{Sampling covariance of the pieces}: the PM total and
$Z$ read the same miss realization, so their seed noise partially cancels in
the subtraction (at $8\times$, $\mathrm{sd}(Z)=.021$ exceeds
$\mathrm{sd}(\hat\sigma_w^2)=.011$); the subtraction reduces variance as
well as bias. (iv) \emph{Plug-in $\hat\sigma_a^2$} propagates one-for-one
($\pm.003$). (v) \emph{Cross-leverage deflation} of the residuals enters PM
and $Z$ with the same sign and largely cancels; second order. (vi) The
\emph{sub-face boundary}: miss below the face scale and the monthly step is
inside $\hat\sigma_w^2$ by convention, a bias only relative to a finer
notion of miss than the method defines. $\hat\sigma_w^2$ remains excluded from the DL weights: with the
miss-adjusted estimator it is now eligible, and its inclusion is an
efficiency study, not a repair; the persistent component alone feeds the
weights until that study is run.

\emph{Development record, kept for its two findings.} The naive diagonal
contrast for $\sigma_w^2$, $(\sum r_i^2-(N-\mathrm{edf}))/\sum w_i$ minus
the persistent and shared terms, is inconsistent: seed spread grows to
$\pm.148$ at $16\times$ because the binomial-scale $(N-\mathrm{edf})$
correction mis-accounts overdispersion absorbed by gate-admitted vertices,
an error that scales with mesh size; the PM trace under working weights is
the repair. And with no persistent component the pre-calibration gate
mass-admitted: overdispersed scores exceed the empirical-null
$\sigma_0\le1.4$ box, selection no trace can price; the dispersion
calibration is that repair, at one measured cost (the stricter gate leaves
slightly more zero-lag miss for the PM total to include: half-mix
$\hat\sigma_w^2$ $0.060$ vs.\ $0.050$ at $16\times$).

Figure~\ref{fig:finalmesh} shows the final-state meshes across regimes; the
center panel is the calibration made visible: pure white pool noise grows
the mesh only $\sim$15\% over the persistent reference (804 vs.\ 700
vertices). A closing battery (three seeds each) confirms no regressions:
chirp attribution $0.117/0.054$, gate insensitivity $0.0448/0.0461$, scale
consistency $0.0476\to0.0452$.

\begin{figure}[h]
\centering
\includegraphics[width=\textwidth]{fig13_final_meshes.png}
\caption{Final-state meshes ($4\times$ data, seed 7, data region). Left:
uniform cosine, pure persistent frailty (1{,}315 faces, 700 vertices).
Center: same target, pure white pool noise (1{,}528 faces, 804 vertices):
the dispersion-calibrated gate admits only $\sim$15\% more vertices where
the uncalibrated gate mass-admitted. Right: chirp (1{,}807 faces, 925
vertices), triangle size tracking local wavelength up the WAC axis.}
\label{fig:finalmesh}
\end{figure}

Autoregressive interpolation ($u_{P,t}=a_P+v_{P,t}$, AR(1) $v$) was
prototyped and shelved: the lag-profile fit identifies the total robustly,
but the permanent/decaying split is horizon-limited (persistence with
$(1-\phi)L\lesssim1$ over an $L$-lag panel is observationally equivalent to
permanence, $L\approx25$--$30$ under cohort attrition) and the lag-zero
extrapolation destabilizes at both $\phi$ extremes. White-plus-permanent
keeps every downstream consumer simple: the credibility posterior is the
static BLUP on $a_P$ with white noise folded into the observation variance.

\section{Theoretical scope and inferential status}\label{sec:scope}

The components of the procedure have different formal status. The table below is intended to keep
those distinctions visible without interrupting every derivation.

\begin{center}
\small
\begin{tabularx}{\textwidth}{p{0.18\textwidth}p{0.18\textwidth}X p{0.20\textwidth}}
\toprule
Component & Status & Main qualification & Principal check \\
\midrule
Candidate score & Local asymptotic approximation & Working WLS neighborhood must be sufficiently converged & Complete-null and score-calibration simulations \\
Local-FDR gate & Conditional empirical-Bayes FDR rule & The candidate family is adaptively generated & Realized FDP and $q$ sweeps \\
EB shrinkage & Working hierarchical model & Depth prior is estimated after selection & Held-out deviance and EB ablations \\
Residual moments & Delta-method and moment identities & Unresolved spatial signal can contaminate pool variance & No-frailty, chirp, and scale experiments \\
WACLS exposure & Exact first-two-moment match & Common conditional event probability within a pool & Heterogeneous-loan-size calibration \\
Pool attribution & Mixed-model/credibility approximation & Requires repeated observations and a stable surface & Persistent-versus-white experiments \\
\bottomrule
\end{tabularx}
\end{center}

\subsection{Conditionality and dataset reuse}\label{sec:conditionality}

One dataset is reused across the empirical null, the gate, the EB
hyperparameters, the variance components, and the pool posteriors. This
section makes the reuse explicit: the estimation order, what each output is
conditional on, where the same data enters an output twice, and what (if
anything) prices that reuse. The intent is that every downstream claim in
this document can be read with its conditioning set attached.

\medskip
\noindent\textbf{Estimation order.} Within a stage: (1) working weights from
the previous stage's $\hat\sigma_a^2$ (binomial at stage 0); (2) refinement
rounds, each comprising candidate scores, the dispersion calibration
$\hat\varphi$, the Lindsey empirical null $(\hat f,\hat\sigma_0,\hat\pi_0)$,
the lfdr prefix gate, EB re-estimation of $\hat\tau_d^2$ over admitted
vertices, and shrunk height updates, with IRLS re-expansion every third round;
(3) on the final mesh, the moment estimators $\hat\sigma_a^2$ (pairwise,
leverage-corrected), the PM total, energy and level, and the BLUP pool
posteriors. Three stages iterate the weights toward their fixed point.

\medskip
\begin{center}
\small
\begin{tabular}{@{}p{2.6cm}p{4.2cm}p{4.6cm}p{3.2cm}@{}}
\toprule
\textbf{Output} & \textbf{Conditional on} & \textbf{Reuse / unpriced selection} & \textbf{Priced by} \\
\midrule
Gate admissions (mesh) & mesh at round; $\hat f,\hat\sigma_0,\hat\pi_0$;
$\hat\varphi$; weights (hence prior-stage $\hat\sigma_a^2$); $\hat\tau_d^2$
via shrunk residuals & candidate $z$'s built from the same residuals that
estimated the null and $\hat\varphi$; hypothesis set itself data-dependent &
empirical null re-centering (partial); FDP measured across seeds: $\approx0.06$ in signal regimes, uncontrolled excursions at the complete null (remark after Proposition~\ref{prop:fdr}) \\
\addlinespace
$\hat\varphi$ (dispersion) & current fit; edf $=$ active-vertex count &
edf produced by the gate that $\hat\varphi$ calibrates; selection absorbs
$>1$ noise unit per admitted vertex & unpriced; direction conservative for
admission \\
\addlinespace
$\hat\tau_d^2$ (EB prior) & admitted set at depth $d$; current weights &
moment over \emph{selected} (truncated) coefficients; feedback into next
round's scores & unpriced (correction queued) \\
\addlinespace
$\hat\sigma_a^2$ & final $\hat g$; mesh; weights & residuals fitted on the
same data; own-pool leverage & analytic correction (\S9.5--9.6), oracle-validated;
plug-in in $\hat\tau,\tilde w$ covered by bootstrap \\
\addlinespace
PM total, $\hat\sigma_w^2$ & final $\hat g$; trace under working weights
$+$ EB ridge; $\hat\sigma_a^2$ subtraction & trace treats the gate-selected
basis as fixed (Efron gap); zero-lag miss inseparable & unpriced; stated as
diagnostic with known confound \\
\addlinespace
Energy, level & final $\hat g$; vertex-count edf & same residuals as all
above; edf at binomial scale & identity checks internal only \\
\addlinespace
BLUP $\hat u_P$ & final $\hat g$; $\hat\sigma_a^2$; weights & pool's own
data in both $\hat g$ and its posterior & shrinkage is the mitigation, not
a correction \\
\addlinespace
Held-out deviance (\texttt{SPLIT}) & even-indexed fit only & none w.r.t.\
surface, gate, and EB (odd pools carry $\approx0$ fit weight) & out-of-sample
by construction \\
\bottomrule
\end{tabular}
\end{center}

\medskip
\noindent\textbf{Reading.} Three circularities are structural: the
weights--$\hat\sigma_a^2$ loop (iterated as the three-stage fixed point,
convergence empirical), the $\hat\varphi$--edf loop (the calibration relies
on the accounting that Remark on selection says is biased; the two errors
push in opposite directions and which wins is regime-dependent), and the
null--gate loop (the null is estimated from the same scores it judges).
Accordingly: FDR statements are per-round Bayes, conditional on the mesh;
$\hat\tau_d^2$ and $\hat\sigma_w^2$ are point estimates with acknowledged
selection bias; $\hat\sigma_a^2$ is the one output with an
oracle-validated correction for its principal reuse channel; and the
held-out deviance is the only fully out-of-sample quantity produced by the
pipeline. Claims elsewhere in the document should be read against this
table.

\section{Numerical evidence}\label{sec:numerics}
\label{sec:oracle}

The review notes, correctly, that no analogue of the
Binev--Dahmen--DeVore/Abramovich rate results is even conjectured here. This
section scopes what such a theorem would have to prove, with an honest
account of what the empirical study can and cannot establish.

The right decomposition separates mesh \emph{placement} from height
\emph{estimation}: at matched vertex count $N$, compare the interpolation
error of the true $f$ on (i) a reference topology built by a noiseless
greedy (projection-gain scoring, no gate, no shrinkage) and (ii) the
topology the gated pipeline selected from \emph{noisy} data; both under
the identical interpolation rule, so that only placement differs
(Figure~\ref{fig:oracle}). The result: the gated topology \emph{beats} the
greedy by a factor of $2$--$3$ at matched $N$ on both surfaces
(uniform: $5.4$ vs.\ $15.8$ weighted RMSE at $N\!\approx\!700$; chirp:
$20$ vs.\ $48$ at $N\!\approx\!925$). Since nothing can beat a genuine
near-optimal oracle, the correct reading is that batched projection-gain
greedy is a \emph{baseline}, not an upper benchmark: it is blind to NVB
closure costs (conformity bisections spend budget its indicator never
prices) and its batch admissions lose greedy fidelity. What the study
establishes is therefore a floor, cleared decisively: the gate's
noise-facing selection places vertices better than a natural noiseless
greedy. The upper benchmark the adaptivity conjecture needs; one-at-a-time
greedy with closure-aware error accounting; is specified and queued; the
conjecture itself would then read: the gated topology's interpolation error
at $N$ vertices is within a constant of that benchmark's, with the constant
degrading gracefully in the FDR target. On the chirp the gated topology's
error--$N$ slope is $\approx-0.7$ against the ideal $N^{-1}$ (two effective
budgets, one seed: indicative only).

Two methodological findings from constructing the study deserve record.
First, an unregularized least-squares fit on an adaptive mesh (the naive
oracle-with-fitting) is \emph{worse at more vertices}: fine vertices with
under-determined stars hold arbitrary heights between data points, invisible
to in-sample metrics and dominant in integrated ones; which is a
back-handed validation of the EB layer, since exactly this pathology is what
the parent-interpolation prior suppresses in the pipeline proper. Second,
the pure-interpolation error of the gated topology ($\approx5$ weighted RMSE
units on the uniform surface) is an order of magnitude below the noisy fit's
error ($\approx60$--$70$): at these information levels the fit is
estimation-noise dominated, so mesh-placement optimality is not the binding
constraint on accuracy; consistent with the cap-sweep finding that
coarser meshes can win.

\begin{figure}[h]
\centering
\includegraphics[width=\textwidth]{fig14_oracle_ratio.png}
\caption{Placement quality at matched $N$: interpolation of the true $f$ on
the gated topology (green) vs.\ a noiseless greedy baseline (black), with
the noisy fit (red) for scale. The gated topology clears the baseline by
$2$--$3\times$; the baseline is not an upper oracle (closure-blind,
batched), so no oracle \emph{ratio} is claimed; the closure-aware
benchmark is queued.}
\label{fig:oracle}
\end{figure}

\subsection{Gate calibration, attribution, and effective exposure}\label{sec:sims}


\begin{figure}[ht]
\centering
\includegraphics[width=0.88\textwidth]{fig21_mesh_scenario_gallery.png}
\caption{Final meshes for matched-support experiments. Face color records mean refinement depth. The complete null remains comparatively coarse; smooth signal induces coherent refinement along support; adding persistent pool variation changes the fitted uncertainty more than the broad geometry; and the chirp spends substantially more depth where the target oscillates rapidly.}
\label{fig:mesh-gallery}
\end{figure}

\begin{figure}[ht]
\centering
\includegraphics[width=0.88\textwidth]{fig23_admitted_vertex_depth.png}
\caption{Locations and depths of active nonstructural vertices in representative smooth and chirp runs. The smooth target distributes refinement broadly along supported cohort bands, whereas the chirp concentrates deep vertices in its fast-scale region. The plot links the statistical gate directly to the geometry it creates.}
\label{fig:vertex-depth}
\end{figure}

\begin{figure}[h]
\centering
\includegraphics[width=\textwidth]{fig1_surface_mesh.png}
\caption{Uniform-cosine target ($\sigma_u=0.30$, $4\times$, 24 pts/face). True surface; fit on its NVB mesh (refinement concentrates on the data support, flat off-support by the zero-mean prior); error masked to support. The quoted precision-weighted RMSE weights each cell by its binomial information.}
\label{fig:surface}
\end{figure}


\begin{figure}[h]
\centering
\includegraphics[width=\textwidth]{fig5_scaling.png}
\caption{Consistency and runtime against data size (single seed per point; wiggles are sampling noise). Runtime is near-linear in $N$; the leverage correction removes the systematic $\approx-9\%$ bias of the naive in-sample $\hat\sigma_u^2$ at every scale; the energy estimator tracks its (shrinking) oracle truth.}
\label{fig:scaling}
\end{figure}


\begin{figure}[h]
\centering
\includegraphics[width=\textwidth]{fig6_sweeps.png}
\caption{The two main dials (single seed). Left: the FDR target $q$ controls mesh complexity against detected miss. Right: naive vs.\ leverage-corrected $\hat\sigma_u^2$ across true frailty levels at $4\times$ data.}
\label{fig:sweeps}
\end{figure}


\begin{figure}[h]
\centering
\includegraphics[width=\textwidth]{fig8_capsweep.png}
\caption{Density-cap sweep on the uniform surface (seeds 7/8/9). Tightening the cap from 12 to 24 pts/face improves precision-weighted surface RMSE on every seed while running $\sim$40\% faster: at $q=0.10$ the gate over-admits vertices that fit frailty-correlated noise, and the cap acts as variance control. $\hat\sigma_u^2$ is insensitive to the cap ($\pm5\%$ of truth throughout) and the no-frailty floor holds. Note the cap binds only at round granularity, so a large admission round can overshoot it (visible in the 48-pts/face spread). This ordering \emph{reverses} on multi-scale truth (Section~\ref{sec:chirp}).}
\label{fig:capsweep}
\end{figure}


\begin{figure}[h]
\centering
\includegraphics[width=\textwidth]{fig9_mesh_compare.png}
\caption{The cap's mechanism, visually: identical data under the 12- and 48-pts/face caps. The fine mesh carries clusters of noise-chasing micro-triangles in the dense band; the coarse mesh tessellates at the target's natural scale and halves the weighted error (most favorable seed shown; see Figure~\ref{fig:capsweep} for the seed spread).}
\label{fig:meshcompare}
\end{figure}



All experiments use the cohort-geometry generator (rate path, origination waves,
attrition) so that support, density, and information heterogeneity match the
application. Selected results from development:

\begin{center}
\begin{tabularx}{\textwidth}{@{}p{0.26\textwidth}X p{0.25\textwidth}@{}}
\toprule
Experiment & Result & Reference value\\
\midrule
Gaussian cosine, uniform sampling & fit error $0.121$, test MSE $1.13$ & noise floor $1.0$\\
Gaussian cosine, cohort geometry & fit error $0.047$, test MSE $1.034$ & noise floor $1.0$\\
Prepayment correction (binomial) & $94\%$ of systematic error removed & coordinate ceiling $\approx 92\%$\\
Frame attachment, cohort truth & Lagrangian rank-1 share $0.997$; $\hat\lambda=0.0196$ & $\lambda=0.02$\\
Frame attachment, location truth & Eulerian rank-1 share $0.989$; $\hat\lambda=0.0250$ & $\lambda=0.02$\\
Overdispersion & $\hat\sigma_u^2=0.087$--$0.088$ & $\sigma_u^2=0.090$\\
Known shift omitted (audit) & $\hat\sigma_u^2=0.201$ & $\sigma_u^2+\beta^2=0.2125$\\
Excess-moment comparison & bias $-0.011$, RMSE $0.038$ (best of five) & truth by construction\\
Double-dip correction & trace $-0.00005$, bootstrap $+0.00016$ & naive $-0.018$; target $0$\\
Oracle pairwise moment (no fit) & bias $\le 1\%$, all configurations & $\sigma_u^2=0.09$\\
$\hat\sigma_u^2$ leverage bracket & in-sample $0.083$, holdout $0.126$ & $0.090$\\
$q$-sweep (0.02--0.5) & miss U-curve, min at $q\in[0.1,0.25]$ & stack $q$-invariant\\
\bottomrule
\end{tabularx}
\end{center}

Notable negative results, retained because they shape the design: empirical-logit
responses cap recovery at $59\%$ (Section~\ref{sec:empiricallogit}); the theoretical
$N(0,1)$ null without the variance correction admits everything and diverges;
unconstrained central matching mid-refinement stops refinement at $\sim 3\%$ of the
needed resolution; pooled (non-slab) $\tau^2$ freezes fit quality under sparse truth;
Eulerian gradient-based estimation of the transport velocity is unidentified because
monthly cross-sections of cohort data are one-dimensional curves (the resolution is to
impose the known kinematics and select the attachment frame by rank-1 separability of
the nodal panel).

\subsection{Balance-weighted outcomes and continuous WACLS exposure}
\label{sec:wacls-sims}

We repeated the main simulation battery with heterogeneous current loan sizes
and balance-weighted prepayments.  Individual balances were lognormal with
coefficient of variation one; the same underlying pools and loan-level
prepayment draws were evaluated once using literal loan counts and once using
the continuous pseudo-counts in \eqref{eq:wacls-counts}.  At this dispersion,
$n_{\mathrm{eff}}$ is roughly one half of the literal loan count, so the
experiment is also a substantial information-reduction stress test.

An isolated calibration with 250 loans and $p=0.20$ shows the reason for the
construction.  The Pearson-residual variance under the ordinary-average-size
proxy rises from $1.06$ at loan-size CV $0.25$ to $1.24$, $2.00$, and $3.20$
at CV $0.5$, $1.0$, and $1.5$.  The WACLS pseudo-count gives respectively
$1.00$, $1.00$, $1.02$, and $1.03$.  The variance correction is therefore
first-order exact, with only the expected finite-pool and Monte Carlo error.

Across 90 matched scenario--seed runs, using continuous rather than rounded
WACLS pseudo-counts changed surface RMSE by $-0.018$ on average (median
$-0.002$), improved RMSE in 53 of 90 runs, and changed the corrected persistent
variance by $-0.003$ on average.  Continuous exposure admitted about seven more
vertices on average, but produced no systematic gate inflation.  In particular,
the complete-null variance floor remained zero and the complete-null surface
RMSE was $0.015$ under both rounded and continuous forms.

The principal scenario results are:
\begin{center}
\begin{tabular}{@{}lrrrr@{}}
\toprule
Scenario & Count RMSE & WACLS RMSE & Count $\hat\sigma_a^2$ & WACLS $\hat\sigma_a^2$\\
\midrule
Uniform surface & 0.434 & 0.370 & 0.073 & 0.074\\
Chirp surface & 0.518 & 0.512 & 0.242 & 0.235\\
No frailty & 0.425 & 0.349 & 0.012 & 0.000\\
Complete null & 0.010 & 0.015 & 0.000 & 0.000\\
\bottomrule
\end{tabular}
\end{center}
The literal-count and WACLS fits become nearly indistinguishable as data
increase.  On the chirp, the WACLS persistent-variance estimate falls from
$0.258$ at $1\times$ to $0.130$ at $4\times$ and $0.100$ at $16\times$, against
truth $0.09$: unresolved fine-scale spatial structure is initially attributed
to the pool component and is reclaimed by the surface with resolution, exactly
as in the count experiments.  The persistent/white decomposition also retains
its ordering under continuous WACLS exposure, and the FDR-target sweep shows no
systematic increase in admissions or empirical FDP.

The learned axis metric remains effective: on the chirp it lowers WACLS surface
RMSE from $0.512$ to $0.240$ and the corrected persistent variance from $0.235$
to $0.113$.  Density caps of 12, 24, and 48 points per face were nonbinding at
the $1\times$ WACLS workload; the gate and information floor stopped first.
The conclusion is operational: use literal counts for count outcomes and
continuous WACLS pseudo-counts for balance outcomes when current-balance WACLS
is available.  Rounding the pseudo-counts is unnecessary and is retained only
for backward comparisons.

\subsection{Multi-scale chirp stress test}\label{sec:chirp}

To test behavior when the target's length scale varies across the state space,
we replace the uniform cosine with a chirp,
$f(x,y)=A\cos\!\big(0.625\,(y+4)^2\big)\cos(2x)$, whose instantaneous frequency
rises linearly from $0$ to $10$ rad/unit along the (fully populated) WAC axis;
total oscillation count matches the uniform surface. Four findings
(Figure~\ref{fig:chirp}).

First, the mesh adapts to local scale: median triangle extent per WAC band
tracks the local wavelength across the support. Second, at fixed data ($4\times$)
the error concentrates at the fine scale: precision-weighted RMSE by frequency
band is $61/75/90$ (slow/mid/fast). Third, and most consequential for the
pool layer: under-resolved fine scale is \emph{attributed to}
$\hat\sigma_u^2$, which rises to $0.12$--$0.16$ against a generative
$\sigma_u^2=0.09$. This is the intended accounting, not a failure: the
estimator's contract is to measure signal left over, and fast-band miss is
pool-persistent (pools ride cohort trajectories at essentially fixed WAC), so
at the current resolution it \emph{is} pool-indexed unexplained signal. The
energy estimator tracks its oracle throughout; only the attribution between
shared-spatial and pool-idiosyncratic moves. Consistent with this reading, the
density-cap ordering flips relative to the uniform surface (finer mesh wins,
$0.122$ at 12 pts/face vs.\ $0.154$ at 24): on multi-scale truth the coarse
cap converts resolvable signal into reported $\hat\sigma_u^2$ rather than
discarding noise, and upward drift of $\hat\sigma_u^2$ under
refinement-starving is the flag that refinement (or a finer cap) still has
signal to claim.
Fourth, consistency restores the estimate: at $16\times$ data with a loose cap
the gate refines to $8{,}407$ vertices, resolves the fast band, and the
leverage-corrected $\hat\sigma_u^2$ returns to $0.0933$ ($+3.7\%$).

\begin{figure}[h]
\centering
\includegraphics[width=\textwidth]{fig12_chirp.png}
\caption{Chirp stress test ($\sigma_u=0.30$, $4\times$ data unless noted).
Left to right: true surface; fit on its adaptive mesh (triangles shrink up the
WAC axis); triangle extent vs.\ local wavelength (medians track the
$\propto\lambda$ reference); weighted RMSE by frequency band, with the
$\hat\sigma^2_u$ recovery $0.127\to0.093$ from $4\times$ to $16\times$.}
\label{fig:chirp}
\end{figure}


\begin{figure}[ht]
\centering
\includegraphics[width=\textwidth]{fig22_chirp_mesh_scaling.png}
\caption{Final chirp meshes at $1\times$, $4\times$, and $16\times$ information under continuous WACLS exposure. The physical WALA--WAC window and target are fixed. Increasing information permits progressively deeper refinement in the high-frequency region, providing a geometric view of the simultaneous fall in surface error and return of the persistent-variance estimate toward its generative value.}
\label{fig:chirp-mesh-scale}
\end{figure}

At $V=8{,}407$ the correction's dense LU factorization became the bottleneck
(183s, $O(V^3)$); since $A_\Lambda$ has $\sim$7 nonzeros per row (P1 elements
plus the 3-point ridge stencil), we replaced it with CSR storage and
Jacobi-preconditioned conjugate gradients for the Hutchinson probe solves,
reproducing the dense results to every reported digit at $O(\mathrm{nnz})$ per
iteration: the correction drops to 1.6s at that size and the pipeline remains
dependency-free.

\subsection{Attribution semantics: what $\hat\sigma_u^2$ means under underfit}
\label{sec:attribution-semantics}

How much of the chirp's $\hat\sigma_u^2$ inflation is recoverable as
\emph{shared} structure? A lag-matched contrast makes the attribution
operational: within-pool pairs sit at month-lags along cohort trajectories, so
we compare them to cross-pool pairs at the same displacement within a WAC band
--- shared spatial miss enters both equally at matched lag, frailty only the
within-pool side, and the difference estimates the idiosyncratic component.
The estimator is exact under the oracle (truth-residual cross moment $=0.0000$
on the uniform surface) and clean in the resolved regime: uniform $4\times$
gives shared $-0.001$ and idiosyncratic $0.079$, reproducing the naive
$\hat\sigma_u^2$. On the chirp the contrast recovers only part of the excess,
and sharpening the matching resolution saturates (shared $0.014/0.021/0.023$
and idiosyncratic $0.140/0.134/0.131$ at $24/48/96$ WAC bands, within-pool
moment $0.154$): the unresolved miss amplitude itself varies below band scale,
so its cross-pool average is recovered while its cross-pool \emph{variance}
remains pool-indexed. This is an identifiability boundary, not an estimator
defect: the shared/idiosyncratic split is defined only relative to a
resolution, and structure below the current resolving power is observationally
equivalent to frailty. The practical diagnostic follows: sweep the refinement
budget (or data) and read the profile of $\hat\sigma_u^2$; a declining
segment is resolvable signal being reclaimed by the surface ($0.154\to0.093$
from $4\times$ to $16\times$ here), and the plateau is the part entitled to
the name frailty.

\subsection{Learned axis metrics, and what the unit-square rescale actually is}
\label{sec:metric}

The rescale of $(\text{WALA},\text{WAC})$ to the unit square is not a numerical
convenience: newest-vertex bisection has no in-plane direction choice (the
refinement edge is dictated by the tree, and the finite similarity-class
structure is what guarantees shape regularity, conforming closure, and the
chord-routing invariant), so the coordinate map is the \emph{only} anisotropy
control, and mapping both axes to $[0,1]$ silently declares one full WALA range
equal to one full WAC range in difficulty at every depth. The algorithmic
version of ``incorporate the relative scales'' is therefore to learn the map:
after the stage-0 pilot fit we estimate density-weighted directional curvature
energies per axis band and equidistribute each axis by
$\Phi'\propto|f''|^{1/2}$ (knot density for $P_1$ error $h^2|f''|$), a
piecewise-linear monotone reparametrization (\texttt{DMESH\_METRIC=1}).

Two findings temper and sharpen the idea. First, on the chirp with a loose
refinement budget the learned metric is nearly immaterial for surface error
(weighted RMSE $70.2\to71.2$): the lfdr gate already performs the
equidistribution itself through refinement depth (Figure~\ref{fig:chirp},
third panel), and the fast band is data-limited, not parametrization-limited.
Second, in the \emph{budget-limited} regime (24 pts/face) the metric
reallocates scarce resolution and acts mainly through the frailty channel:
corrected $\hat\sigma_u^2$ falls from $0.162$ to $0.124$ (against truth
$0.09$) with $10\%$ fewer vertices, because thinner fast-band strips make the
residual miss vary within pools rather than ride them, cutting the
pool-persistent leakage into the pairwise moment. The uniform-surface results
are unchanged under the learned map, and the no-frailty floor rises mildly
($0.0003\to0.0010$; the pilot fit's noise curvature warps the map when there
is no signal), so the option defaults off. Genuinely local anisotropy
(directional ratios varying point to point) is out of reach of any separable
map and of isotropic NVB; the general tool is metric-field remeshing or
dimension-adaptive tensor bases, which we note as the larger lift. Single-seed
numbers throughout this subsection.


\subsection{Implementation cost}
\begin{figure}[h]
\centering
\includegraphics[width=0.75\textwidth]{fig7_optimization.png}
\caption{Implementation cost at the $8\times$ workload (46k pool-months), log scale: Python baseline through the C++ port, stage-0 cap, Hutchinson diag($S$), container refactor, and heap disable. The subsequent sparse-CG solve (Section~\ref{sec:chirp}) removes the remaining $O(V^3)$ term in the correction.}
\label{fig:speed}
\end{figure}
\section{Real-data check: Ginnie Mae pools and the EB ablation}\label{sec:realdata}

We ran the full C++ pipeline on the Ginnie Mae single-family pool portfolio
disclosure for April 2026 (\texttt{monthlySFPS\_202604}): 120{,}595 SF pools
covering 47M loans after basic filters, including the multi-issuer mega-pools
(up to $\sim$150k loans). A single monthly snapshot carries no monthly
prepayment outcome, but original loan count is recoverable as
$n_0=\mathrm{round}(\text{OrigAggregate}/\text{AOLS})$, making
$(n_0-n_{\text{now}}\mid n_0)$ a binomial \emph{cumulative termination}
outcome per pool.  Here $n_{\text{now}}$ is the disclosed current number of
loans, not Pool UPB divided by an average loan size: the latter falls under
scheduled amortization even without termination.  A balance-weighted version
would instead require current-balance WACLS and the continuous exposure in
\eqref{eq:wacls-counts}, which is not disclosed in the pool-level file. The gated mesh fits the cumulative termination surface over
$(\text{WALA},\text{WAC})$ from a constant-offset start in a few seconds
(3{,}103 vertices under the final dispersion-adjusted gate), with the
one-step re-expansion carrying the fit across fitted probabilities spanning
$0.03$--$0.99$; the level statistic lands at $+0.013$. The dispersion
calibration of Section~\ref{sec:poolstate} improved the real-data fit
materially: residual overdispersion beyond binomial (the energy read) fell
from $0.25$ to $0.16$, the uncalibrated gate having over-admitted on real
pools exactly as in the pure-white simulation.
With one observation per pool the pairwise moment has no within-pool pairs, so
$\sigma_u^2$ is structurally unidentifiable and is held at its floor; a
two-month panel (loan-count SMM) identifies it.


\begin{figure}[h]
\centering
\includegraphics[width=\textwidth]{fig10_real_data.png}
\caption{The real-data fit. Fitted cumulative termination probability over $(\text{WALA},\text{WAC})$ with color fading where data is thin; the same surface on its adaptive mesh; and the seasoning ramp by WAC tercile (dots observed, count-weighted; lines fitted), where the coupon separation is the refinance-incentive signal read directly from pools.}
\label{fig:realdata}
\end{figure}

The real data gives the shrinkage layer a demanding test: the median GNMA~II
custom pool contains 13 loans at origination, so thousands of mesh vertices
are supported by a handful of tiny pools. An ablation isolates the effect:
fitting on even-indexed pools and scoring held-out odd-indexed pools, with
empirical-Bayes shrinkage on versus off (raw local WLS heights).

\begin{figure}[h]
\centering
\includegraphics[width=\textwidth]{fig11_eb_ablation.png}
\caption{EB ablation on held-out real pools. Removing shrinkage raises
held-out deviance by 60\% (13.4 to 21.5 per pool), nearly triples Pearson
$X^2$ per unit information (0.35 to 0.95), and degrades the level calibration
tenfold ($-0.006$ to $-0.062$). Without shrinkage the candidate $z$-scores are
also wilder, so the gate over-admits (3{,}290 vs.\ 2{,}854 vertices),
compounding the overfit. The only quantity EB costs is fit time (2.5s vs.\
1.1s).}
\label{fig:ebablation}
\end{figure}

The ablation confirms on real data the small-cell noise mechanism that
motivated the hierarchical prior: unshrunk heights on thin support are
essentially noise, and the depth-indexed EB precision $1/\tau_d^2$ removes it
at the cost of a factor of two in fit time.


\section{Updated controlled benchmark and deep-refinement diagnostics}\label{sec:updated2026}

This section records the July 2026 revisions to the simulation design and the subsequent
Ginnie Mae diagnostics.  These experiments sharpen two distinct questions that were
confounded in the earlier development runs: whether the mesh places resolution correctly
when the sampling support is benign, and whether hierarchical empirical-Bayes shrinkage
remains sufficiently strong after many successive refinements in real data.

\subsection{Uniform-support benchmark with increasing spatial frequency}

The default controlled simulation now removes the cohort-support artifact.  WAC is drawn
independently over its full range, every simulated pool is observed at WALA 1--60, and
there is no attrition.  The displayed horizontal coordinate is
\[
 x=\mathrm{WALA}/12\in[1/12,5].
\]
The target remains the two-scale exponential-cosine surface, but its normalized horizontal
coordinate $z\in[0,1]$ is warped according to
\[
 z_{\mathrm{warp}}=(0.2+0.8z)z.
\]
Thus the local horizontal frequency increases smoothly toward larger WALA while the data
remain essentially uniform over the rectangle.  This design is intended as a geometric
identification experiment rather than a realistic mortgage-vintage model.

In the representative seed-10 run, the fitted surface had grid correlation approximately
$0.992$ with the known truth and grid RMSE approximately $10.6$ centi-log-odds.  The mesh
visibly concentrated resolution where the warped target oscillated most quickly.  The
controlled experiment therefore indicates that the refinement machinery can track
spatially heterogeneous frequency when data support is not itself wedge-shaped.

\begin{figure}[ht]
\centering
\includegraphics[width=\textwidth]{fig24_uniform_warped_benchmark.png}
\caption{Controlled uniform-support benchmark with increasing horizontal frequency.  The
panels show the exact truth, adaptive-mesh estimate, estimation error, and the final mesh.
The data cover the full WALA--WAC rectangle with zero attrition, separating geometric
adaptation from sampling-support artifacts.}
\label{fig:uniform-warped}
\end{figure}

\subsection{June 2026 Ginnie Mae fit on the log-odds scale}

The same estimator was applied to June 2026 Ginnie Mae single-family pool data.  Because
there is no known ground-truth correction surface, real-data diagnostics are reported on
the log-odds scale using information-weighted empirical logits and held-out residuals.  The
full fit used 371,089 pools, 8,601 vertices, and 17,049 triangles.  The resulting geometry
is substantially more irregular than in the controlled benchmark because the empirical
WALA--WAC support and information density are highly nonuniform.

\begin{figure}[ht]
\centering
\includegraphics[width=\textwidth]{fig25_gnma_logodds.png}
\caption{June 2026 Ginnie Mae fit on the log-odds scale.  The panels show empirical log
odds, fitted log odds, held-out empirical-minus-fitted log odds, and the adaptive mesh.
Empirical logits use a Jeffreys $0.5$ correction to avoid infinities in saturated cells.}
\label{fig:gnma-logodds}
\end{figure}

\subsection{Evidence of undershrinkage after deep refinement}

The real-data mesh contains clusters in which triangle area becomes extremely small after
many successive refinements.  Several diagnostics indicate that at least some of these
clusters retain too much fine-scale variation.  First, held-out error is higher in the
finest triangle-area decile than at several intermediate resolutions.  Second, a
cross-fitted multilevel diagnostic compares the variation implied by the sequence of
hierarchical prior variances with independently realized variation over the final $k$
levels of the hierarchy.  For a suspicious central cluster the realized-to-implied
variance ratio was approximately
\[
 0.90,\ 0.94,\ 1.40,\ 0.35
\]
for one through four levels respectively.  Thus one-level behavior is not obviously
miscalibrated, while four consecutive levels collectively imply roughly three times the
variation supported by independent data.

Depthwise prior standard deviations in the same experiment became very large at
intermediate-deep levels, reaching roughly $1.9$--$2.1$ log-odds units at depths 6--7.
Turning on the existing admission-truncation correction improved held-out fit but drove
most deep prior scales to the lower floor, suggesting that the uncorrected estimator is
too permissive whereas the full correction is too aggressive.

\begin{figure}[ht]
\centering
\includegraphics[width=0.94\textwidth]{fig26_multilevel_tau.png}
\caption{Multilevel hierarchical-variance diagnostic on a stratified Ginnie Mae sample.
Left: cross-fit realized variance divided by the variance implied by the estimated
hierarchical prior over multiple levels.  Right: estimated prior standard deviation by
depth under the ordinary and admission-truncation-corrected moments.  The ordinary
estimator permits excessive cumulative variation over several consecutive levels, while
the full truncation correction collapses deep scales toward the imposed floor.}
\label{fig:multilevel-tau}
\end{figure}

A direct manual ablation confirms that the smallest-triangle clusters contain removable
fine-scale variation.  Holding the geometry fixed, a local Laplacian smoothing step was
applied only to the finest 10\% of vertices.  On a stratified 80,024-row sample, the best
manual setting reduced held-out binomial deviance per pool from $19.29$ to $18.08$ overall.
Within the support directly affected by the smoothing, deviance fell from $53.50$ to
$23.64$.  The improvement is heterogeneous: some central regions did not benefit from the
same blanket smoothing rule.  The conclusion is therefore not that all deep refinement is
spurious, but that the current shrinkage can be insufficient in selected runaway branches.

\begin{figure}[ht]
\centering
\includegraphics[width=0.86\textwidth]{fig27_manual_smoothing_ablation.png}
\caption{Held-out deviance under manual smoothing of only the finest mesh vertices.  The
ablation demonstrates that additional smoothing in selected deep clusters can materially
improve predictive likelihood, motivating a more selective empirical-Bayes calibration of
fine-scale branches rather than a global increase in regularization.}
\label{fig:manual-smoothing}
\end{figure}

\subsection{Focused diagnostic: WALA 100--150 and WAC 3.5--4.5}

A focused window at WALA 100--150 and WAC 3.5--4.5 contains 5,075 observations, 309 active
vertices, 123 gate-admitted vertices, and 678 intersecting triangles.  Its held-out
log-odds RMSE is $1.05$ with mean empirical-minus-fitted log-odds error $-0.18$.  The
hierarchical surplus magnitude does not decay with depth: median absolute surplus is about
$0.48$ at depth 7, $0.72$ at depth 8, $0.96$ at depth 9, and $0.87$ at depth 10.  All five
local depth-10 vertices are admitted.  This is consistent with a refinement branch that
continues to carry large coefficients after the mesh has become very fine.
These observations are resolved by the certified benchmarks of
Section~\ref{sec:certnull}: the deep admission \emph{count} in such
windows is reproduced on data whose truth is provably flat there (an
admission-side artifact of the conditional null under real-design density
and overdispersion), while the anomalous \emph{amplitude} decomposes into
mega-pool state masquerading as surface plus a residual genuinely sharp
component; the composite gate of Section~\ref{sec:nulls} removes the
manufactured population ($886\to2$ on the certified null) and the manual
smoothing ablation below is the crude form of the post-selection
re-shrinkage pass now implemented.

\begin{figure}[ht]
\centering
\includegraphics[width=\textwidth]{fig28_gnma_region_zoom.png}
\caption{Focused Ginnie Mae diagnostic for WALA 100--150 and WAC 3.5--4.5.  The panels
show empirical log odds, fitted log odds with the local mesh, held-out residual structure,
and the mesh with the finest local triangles emphasized.}
\label{fig:gnma-region-zoom}
\end{figure}

\begin{figure}[ht]
\centering
\includegraphics[width=0.88\textwidth]{fig29_gnma_region_depth.png}
\caption{Hierarchy diagnostics in the WALA 100--150, WAC 3.5--4.5 window.  The number of
admitted vertices remains substantial at deep levels and median surplus magnitude rises at
depths 8--10 instead of dying away.}
\label{fig:gnma-region-depth}
\end{figure}

\paragraph{Current recommendation.}
The recommended estimator remains the original local-FDR adaptive mesh with hierarchical
empirical-Bayes shrinkage.  A first broad local-branch EB model was tested and rejected:
it improved some suspicious clusters but overshrank healthy regions and worsened aggregate
held-out performance.  The open methodological problem is therefore to regularize only
branches whose \emph{multilevel} implied variation is unsupported, rather than to impose a
global fine-mesh penalty.  The controlled uniform-support benchmark and the real-data deep
branch diagnostics are both included in the accompanying code package for reproducibility.


\section{Certified null benchmarks}\label{sec:certnull}
The certified-window instruments separate predictive fit from false geometric discovery. They use
the exact real design but replace the truth inside a declared window by a function whose deep
hierarchical surpluses are negligible. A deep admission in the interior is therefore false by
construction. This section reports two different statuses: the legacy gate has passed these
instruments; transfer of the same claims to the exact-solver gate has not.

An earlier matched-null claim was retracted on a ground worth stating as a result: matched fidelity
and sharp-null-everywhere are incompatible. The real termination ramp carries depth-8-scale
curvature near WALA $0$, so no benchmark can match the ramp and be null there. The standing
instruments instead certify
$W=\mathrm{WALA}[80,200]\times\mathrm{WAC}[3,6]$ on the exact April 2026 design.

\subsection*{Legacy certified results}
\textbf{v3.} Truth is a realistic nine-parameter log-age polynomial everywhere except $W$, where it
is that polynomial's own quadratic fit with a cosine taper. The maximum true hierarchical surplus
in the window interior is $1.65/0.77/0.36$ centi-log-odds at depths $8/9/10$, two orders of magnitude
below admission scales. The legacy baseline gate produced $886$ certified-false deep faces; the
analytic composite gate produced $2$.

\textbf{v4.} The mean is the same, but the noise includes the measured size-dependent pool variance,
small-pool heavy tails, real issuer assignments with $0.09$ issuer variance, diversified
multi-issuer pools, and a shipped known $n_{\mathrm{eff}}$. The legacy baseline produced $997$
certified-false faces; the analytic composite gate produced $7$ and the permutation gate $6$. The
$2\to7$ change is the expected issuer-clustering attack: co-located pools sharing an issuer draw
replicate like surface unless effective replication is priced over issuers rather than pools.

\subsection*{Exact-score transfer experiment B3: failed}
The B3 transfer froze the signed-root score definition \eqref{eq:exactgain} and attempted to
standardize it by a candidate-specific sandwich factor,
\begin{equation}\label{eq:b3score}
z_v^{B3}=\frac{r_v}{(\widehat{\mathrm{sd}}_{\delta,v}/100)\sqrt{I_v}},
\end{equation}
with the remaining family-level location and scale estimated by central matching. On an independent
complete-null simulation, the first exact-score rounds were close to unit scale (stage 1 mean
$0.013$, sd $1.001$; stage 2 mean $0.012$, sd $0.992$), and the run ended with four false
admissions. Prediction on the original and eruption-seed splits also remained stable.

The certified windows nevertheless rejected the proposal. An initial wiring omission set the
pool-noise variance to zero and reproduced the old failure almost exactly: v3 had $992$ and v4 had
$982$ certified-false faces. After correctly pricing the declared composite variance $0.35$, v3 fell
to $80$ under the standard stage-0 policy and $75$ when stage 0 also used the exact score, still far
from the predeclared target near $2$. The corrected v4 run generated a large adaptive workload and
did not complete within the run budget; its trajectory was not approaching the required $3$--$7$
range. B3 is therefore not accepted, and no legacy certification claim transfers to the exact gate.

The failure is informative. The exact one-dimensional objective and its signed gain pass their local
invariants, but a single global empirical-null/prefix rule can admit null candidates inside $W$ as
passengers behind strong candidates elsewhere or can pool together fit states with different residual
laws. The next calibration must be local to the candidate's held-out evidence or otherwise condition
on the adaptive state.

\subsection*{Conformalized held-out certification: proposed}
A conformal construction is possible, but a naive pointwise conformal interval is not enough. Standard
split-conformal coverage is marginal for a future observation; it does not by itself control an
adaptively generated family of local surface hypotheses. A valid version must separate hypothesis
generation from calibration and preserve the pool/issuer dependence structure.

One concrete design is:
\begin{enumerate}[leftmargin=2em,itemsep=2pt]
\item Split observations by exchangeable pool or issuer blocks into a \emph{construction} sample and
a \emph{calibration} sample. Build the current mesh, fit the null surface, and freeze the candidate
set using only the construction sample.
\item For candidate $v$, estimate its direction and amplitude $\hat a_v$ on the construction sample.
On each calibration block compute the held-out directional log-likelihood increment
\[
A_{bv}=\sum_{i\in b}c_i\left\{\ell_i(\eta_i+X^H_{iv}\hat a_v)-\ell_i(\eta_i)\right\}.
\]
This prevents the candidate from choosing its direction on the same outcomes used to certify it.
\item Within prespecified exchangeability strata (pool size, fitted hazard, issuer structure, and
possibly stage/depth), permute residual bundles among geometric slots or use leave-one-block-out
nonconformity ranks. The conformal/permutation p-value is
\[
p_v=\frac{1+\#\{b:T^{(b)}_v\ge T^{\mathrm{obs}}_v\}}{B+1},
\qquad T_v=\sum_b A_{bv}.
\]
Equivalently, invert the test to obtain a conformal interval for the held-out gain or coefficient and
admit only when the interval excludes zero in the construction-sample direction.
\item Apply a multiplicity rule to the frozen family. Ordinary BH requires its usual dependence
conditions; blockwise permutation p-values or e-values with e-BH are safer when neighboring
candidates overlap. Cross-fitting can rotate construction and calibration folds to recover power
without reusing a candidate's own calibration outcome for its direction.
\end{enumerate}

This is best viewed as a conformalized version of the existing permutation null, with the important
addition that evidence is evaluated out of sample. It offers finite-sample validity under the stated
block exchangeability, but not assumption-free validity: calendar shocks, nonexchangeable issuer
formation, and adaptation after seeing calibration outcomes would break the guarantee. The v3/v4
windows provide the predeclared falsification suite for this proposal before any real-data adoption.


\subsection*{Complete-null comparison and sequential diagnosis}
A later audit separated the certified smooth-window question from a genuine complete null.  Both
complete-null designs reuse the Ginnie covariates and exposures but set the spatial correction to a
constant.  In the clean version the only randomness is binomial; in the matched heavy-tail version
pool effects are drawn from the same four size-stratified two-component law supplied to the
histogram gate.  Ten frozen seeds were run for each design.  The clean null stopped at the root mesh
in every seed: no admitted free coefficients, no depth-$\geq 8$ faces, and fitted spatial rms
$0.37$ centi-log-odds.  The matched heavy-tail null admitted an average of $2.4$ free coefficients
(median $1$, range $0$--$11$), admitted at least one coefficient in $6/10$ runs, produced an average
of $9.4$ depth-$\geq8$ faces, and left $8.09$ centi-log-odds of fitted spatial rms.  Thus ordinary
binomial null testing is not the difficulty; the remaining failure is specific to lumpy pool
heterogeneity and the adaptive reuse of its realized spatial pattern.

The certified v3 instrument is not a complete null.  It contains a legitimate quadratic trend in the
window and certifies only that deeper nonquadratic surpluses are negligible.  Under the deterministic
histogram gate, v3 produced $259$ admitted free coefficients overall and $14$ coefficient centers in
the window, all at hierarchical depths $4$--$5$; none was a depth-$\geq8$ free coefficient.  The
same run produced $35$ area-depth-$\geq8$ faces and approximately $5.65$ centi-log-odds of
nonquadratic fitted texture after removing the best quadratic.  The historical face count therefore
mixes legitimate low-order fitting, NVB closure, and false fine-scale texture.  It remains a useful
end-to-end geometry diagnostic, but it is not by itself candidate FDR after geometric existence and
statistical freedom have been separated.

The complete-null failure was then traced round by round.  The heavy-tail runs execute many adaptive
candidate-family tests; problematic seeds had $18$--$24$ rounds, with six to nine rounds containing
at least one rejection.  At the initial root family, only two of ten seeds rejected anything (six
candidates in one seed and one in another), which is not by itself a precise rejection of a nominal
$q=0.10$ all-null family with only ten repetitions.  Once a coarse false split is admitted, however,
the same realized pool anomaly changes the fit, generates descendants, and receives repeated
opportunities to be selected.  The final admission counts $11,0,5,0,2,0,1,0,1,4$ show this branching
pattern.  The failures are not confined to one-pool stars: among the $24$ admitted coefficients the
median variance-effective pool count was about $55$ and the median fourth-moment effective count
about $21$.  Increasing explicit mixture enumeration from eight to fourteen dominant contributors
left the problematic admission counts unchanged, so truncation of the deterministic histogram
convolution is not the leading error.

The current diagnosis is consequently two-part.  First-round histogram p-values are mildly
anti-conservative in their extreme tail under the lumpy law.  More importantly, per-round BH at a
fresh $q$ does not provide whole-run control for a data-adaptive branching dictionary.  The relevant
null distribution after the first split is selection-conditioned,
\[
 \mathcal L(U_v\mid\text{the current mesh was selected from the same realization}),
\]
not the marginal histogram law used at the root.  The next falsification suite therefore has three
arms: (i) a one-family heavy-tail null over at least $100$ frozen seeds; (ii) a frozen-dictionary
repeated-refit experiment; and (iii) the full adaptive run.  Their differences isolate marginal tail
calibration, refitting reuse, and descendant generation.  A plausible repair is branchwise or
sequential error spending, in which a root branch receives a fixed budget and descendants inherit
subdivisions of that budget rather than receiving a fresh $q=0.10$ family at every generation.  This
proposal is not yet implemented or certified.

\section{Measured noise physics and formation information}\label{sec:physics}

\textbf{Effective exposure has a ceiling.} A pool's information about the
surface saturates at $1/\sigma^2$: $n_{\mathrm{eff}} = n/(1+npq\sigma^2)$,
about $18$ loans at typical rates and $\sigma^2=0.35$. The April file's
$47.6$M loans are $4.2$M effective loans for surface estimation, and its
$32$ mega pools are ${\approx}18$ effective loans each. Fitting, testing,
or scoring at nominal $n$ where $n_{\mathrm{eff}}$ belongs is the single
mechanism behind the cluster pathology, the deviance tail, and the null
miscalibration.

\textbf{The mixing law is not Gaussian.} The per-size instrument
(face-level cross-pool moments remove local surface bias; falsified clean
on the null benchmark) reads $\sigma^2(n)$ near-flat at $0.26$--$0.47$.
The shape test against the fitted binomial-Gaussian mixture rejects:
standardized-residual tails run $12.5\times$ the mixture at $|r|>6$ and
$77\times$ at $|r|>10$ (excess kurtosis $67.7$ against $3.3$), while pools
with $n\ge10^4$ are lighter-tailed than modeled. The mixing law is
size-dependent in shape; the permutation null of
Section~\ref{sec:nulls} prices this without modeling it.

\textbf{Pool formation is an insider partition.} Within-cell termination
rises $+0.11$ per e-fold of loan count and $+0.88$ per e-fold of average
loan size; the loan-size loading sign-flips with seasoning ($-0.21$ young,
$+0.79$ seasoned: credit-event and refinancing channels load on balance
with opposite signs). Issuer identity alone explains $18.9\%$ of residual
within-cell pool variance (issuer sd $0.30$ log-odds), and multi-issuer
pools, which have no single grouping agent, carry $5.4\times$ less
residual variance than single-issuer pools. Grouping metadata forecasts
the pool effect; in an overlay deployment the base model is assumed to own
borrower-side balance effects, and formation and servicer effects are the
overlay's cargo.

\section{Evaluation metrics and the reporting standard}\label{sec:metrics}

Held-out deviance per pool is verified correct and tail-fragile: the top
$100$ of $60{,}297$ held-out pools carry $16.5$ of the baseline's $24.54$,
and the $19$ mega pools average $9{,}475$ each (against ${\approx}1{,}359$,
near their irreducible floor, under the composite gate). It is a
$19$-pool story wearing a $60{,}000$-pool denominator. The headline metric
is therefore the held-out marginal likelihood with the pool effect
integrated out (Gauss--Hermite at fixed $\sigma_{\mathrm{pool}}$), which
scores each pool at $n_{\mathrm{eff}}$ and bounds its influence. The
outlier tail itself was forensically informative: $95\%$ same-sign errors
traced to the within-cell size effect above, an omitted formation
covariate rather than noise. The reporting standard is the triple
(marginal NLL, information-weighted deviance in the differential-refinement
zones, certified-null deep count); a configuration dominates only if it
wins all three, no single-snapshot configuration does, and the panel or
loan-level success criterion is exactly to match the conditional bound on
the first while matching the composite bound on the second and third.


\section{The loan-level formulation}\label{sec:loanlevel}

Balance-dependent refinancing removes exchangeability of loans within a
pool: the future data model treats each loan as its own trial, sharing
only its pool's effect,
\[
Y_{Pj}\sim\mathrm{Bernoulli}\bigl(\operatorname{logistic}(\eta_0(b_{Pj},z_P)
+ g(z_P) + u_P)\bigr),
\]
with $b_{Pj}$ the known loan balance. Seven consequences, each validated
by simulation (\texttt{loanlevel\_sim\_issues.py}; $\beta=0.5$ per
log-balance, $\sigma^2=0.35$).

(1) The conditional law is Poisson-binomial and under-dispersed:
$\operatorname{Var}(k\mid u)=\sum p_jq_j$ (simulated $87.9$, formula
$88.0$) against the binomial $n\bar p\bar q = 91.7$; the binomial noise
floor is an over-estimate. (2) Aggregation does not commute with the
link: at fixed mean $\eta$ the expected rate falls with balance
dispersion ($0.646\to0.625$ across dispersion $0\to1.8$), and pool-level
regression on mean balance recovers $0.42$ for a true $0.5$. The
estimand must move to the loan level; pools are known weighted integrals
of it. (3) Information is concave: $\sum p_jq_j \le n\bar p\bar q$, so
nominal pool information overstates loan-sum information. (4) The pool
effect is a one-factor structure with loadings $p_jq_j$: near-saturated
loans barely reflect their pool's effect; pairwise covariances follow
$p_jq_j\,p_{j'}q_{j'}\sigma^2$; loading-correct moment estimation
recovers $\sigma^2 = 0.320$ where naive $npq$ weighting gives $0.264$
(truth $0.350$). (5) The residual moment bias is an expansion-point
artifact: anchoring the delta method at each pool's posterior mode
$\hat u_P$ rather than the prior shrinks the error base $22\times$
(posterior variance $0.016$ against $\sigma^2=0.35$) and the
fixed-point estimator recovers $0.347$. The same-data ladder closes the
question of higher orders: mode Laplace $0.36407$, fully exponential
second order $0.36488$, exact mode-centered quadrature $0.36489$; the
correction closes the deterministic gap to $5\times10^{-6}$ while the
exact estimate sits $0.7$ sampling deviations from truth. The
approximation is not the binding error. Production recipe: Laplace above
a curvature threshold, mode-centered GH below, second order as
cross-check.

(6) Identification is the headline. Within-pool balance contrasts
difference out $u_P$: the fixed-effects fit recovers
$\hat\beta = 0.497$ even under informative formation
($\operatorname{corr}(u,\text{balance})=0.6$), where pooled estimation
reads $0.651$: a $30\%$ upward bias, pool effects laundered into the
balance curve. For specified pools, balance-homogeneous by construction
(within-pool sd $0.05$), the FE standard error explodes $18\times$ and
pooled reads $0.812$: the confounding concentrates exactly where
formation is most informative, so balance coefficients in spec-pool
strata are identified only across pools and must be flagged as
model-dependent. (7) Estimation: the pool-level Laplace mixed likelihood
(one 1D log-concave integral per pool, $u$ integrated against the
Bernoulli product) jointly recovers $(\hat\beta, \hat\sigma^2, \hat a_0)
= (0.503, 0.365, 0.604)$ on $120$k loans in $1{,}500$ pools; this is the
estimation engine for the loan-level mesh, with issues (1), (3), and (4)
priced automatically. The permutation null generalizes by permuting
pools carrying their loan bundles intact, conditioning on realized
compositions.

The architectural consequences: the mesh, gate, and lfdr machinery are
unchanged; the working linearization per loan sits at
$\eta_0 + g + \hat u_P$ (posterior anchoring as the linearization
discipline); $K_{\mathrm{eff}}$ counts pools of loans; and within-pool
balance variation finally splits the single-snapshot bundle, since pool
idiosyncrasy is constant across a pool's loans while sub-face surface
varies with their coordinates.


\section{Tail shape: the kurtosis program}\label{sec:kurtosis}

Status: all items measured on controlled benchmarks unless labeled
otherwise. The pool-effect law's shape, not just its variance, is a
load-bearing input. The measured excess kurtosis of standardized real
residuals is $67.7$ pooled across sizes (within-stratum recalibration
below: the pooled figure mixes scales); a Gaussian-mixed binomial compound produces
essentially zero at any $\sigma$ (measured: $-0.05$ across
$\sigma^2 = 0.12$ to $2.0$), so the kurtosis is a direct measurement of
the non-normality of $u$. The controlled instrument is a
variance-preserving lump mixture: with probability $\pi(n)$ the pool
effect is scaled by $L$, normalized so $\mathrm{Var}(u)$ is unchanged;
$(\pi, L) = (0.035, 17)$ reproduces the real kurtosis. At that setting
$91\%$ of the effect variance is carried by $3.5\%$ of pools.

\textbf{Translation to the null.} For a candidate with effective
replication $K_{\mathrm{eff}}$, the null statistic's excess kurtosis is
$\kappa_4(u)/K_{\mathrm{eff}}$ (simulation matches across two decades).
Deep candidates on the real design run at median $K_{\mathrm{eff}}
\approx 2$ to $3$ (measured via a score diagnostic), so effect kurtosis
arrives at deep admission decisions nearly undiluted:
Gaussian-equivalent three-sigma cuts are $z \approx 7.5 / 6.2 / 5.1$ at
$K_{\mathrm{eff}} = 2/4/8$. A pipeline experiment with the effect
kurtosis dialed to $\{0, 10, 33, 69\}$ confirms the in-window null
score kurtosis climbs accordingly (e.g.\ depth-6: $6.7 \to 26.9$),
with the caveat that candidate populations under default dynamics are
selection-conditioned; a forced-uniform-refinement rerun is queued to
isolate the marginal law.

\textbf{What kurtosis does not damage.} On a clean benchmark (cosine
product signal, uniform locations, $500$k pools, kurtosis dialed at
fixed variance) the surface fit is unaffected to slightly improved:
rms error $13.99 \to 12.30$ centi-logits from kurtosis $0$ to $69$, and
held-out deviance in excess of the oracle is constant at $0.54$ to
$0.58$. Lumpy noise concentrates the same variance in fewer pools; the
information-weighted average is dominated by the quieter majority
wherever replication is high. Damage is confined to low-$K_{\mathrm{eff}}$
decisions, certification counts, and per-pool posterior tails.

\textbf{The posterior transform.} Posterior effect estimates carry
kurtosis $\lambda^2 \kappa_u$ with $\lambda = \sigma^2/(\sigma^2 + v)$
the pool's reliability; exact at moderate kurtosis, attenuated at the
extreme by outcome saturation. The law inverts to a measurement tool:
divide observed posterior kurtosis by $\lambda^2$ per size class.

\textbf{Estimator diagnosis} (Figure~\ref{fig:sigmadiag}). At kurtosis
$69$, truth $\sigma^2 = 0.35$: robust scale reads $0.065$ (the quiet
core), Gaussian ML $0.108$ (a misspecified-model compromise), the
moment estimator $0.292$ (unbiased in principle; the deficit is outcome
saturation). Counts are bounded, so $k = n$ censors lump magnitudes;
the lump law stays identified through saturation frequencies, which is
what the repair uses. Details and the refuted intermediate hypotheses
are in the lab notes.

\textbf{The repair, validated.} An EM fixed point with the lump-mixture
prior and exact binomial likelihood (mode-centered Gauss-Hermite,
deterministic, no sampling, no logits) recovers all parameters at every
kurtosis level: at $69$, $\hat\pi = 0.0353$ (truth $0.035$),
$\hat\sigma_L^2 = 8.96$ (truth $9.13$), total variance $0.361$ (truth
$0.35$), implied kurtosis $62$; at kurtosis $0$ it degenerates
gracefully (near-equal components, implied kurtosis $0.1$, no phantom
tails). Residual $+4\%$ variance bias is surface-error leakage,
quantitatively matched to the measured fit rms. One estimation pass
supplies the total $\sigma^2$ for shrinkage, $(\pi, L)$ for the
four-moment or saddlepoint gate, and tail-aware posteriors that shrink
identified lumps less (mixture-posterior kurtosis $58$ versus $21$
under the Gaussian prior; correlation with truth $0.94$). Proposed and
queued: wire these into the admission gate; acceptance test is the
full-strength tail benchmark with pass criterion $53 \to$ single
digits.


\textbf{Real-data measurement} (Figure~\ref{fig:ginniemix}). Running the
mixture fixed point on the April 2026 file per size stratum: totals
$0.851 / 0.514 / 0.494 / 0.391$ for $n$ in $25$--$128$ / $128$--$512$ /
$512$--$4096$ / $4096{+}$, with wide-component weights $0.11$ to $0.27$
and scale ratios $5$--$6\times$: a broad two-population structure
rather than rare monsters. Within-stratum residual excess kurtosis is
$1.8 / 1.9 / 2.8 / 20.1$ with $P(|t|>4)$ of $12/20/27/194\times$
Gaussian: the pooled kurtosis $67.7$ was substantially a cross-size
scale-mixture artifact, and the heavy-tail problem concentrates at mega
pools, which are also the $K_{\mathrm{eff}}$-collapsing,
metric-dominating objects. Identified lump pools are one-sided with
opposite signs by size: small-pool lumps $86\%$ slow (mean $-1.35$
logits; an $n_0$-rounding data-artifact channel is queued for checking
before these are treated as physical), mega-pool lumps $73\%$ fast
(buyout-like). Consequences, proposed: per-stratum gate variances from
the fixed-point table (the hand-set $0.35$ underprices small-$n$
strata); tail-aware nulls needed most at mega-star candidates; the
full-strength adversarial benchmark retained as stress test but a
realistic-tails acceptance benchmark (mild small-pool tails, heavy
one-sided mega tails) replaces it as the pass criterion.

\begin{figure}[t]
\centering
\includegraphics[width=0.98\textwidth]{ginnie_mixture_diagnosis.png}
\caption{The real pool-effect law by size stratum. (a) Central shape.
(b) Tail survival curves: mega pools carry the tail weight. (c)
Posterior effect law. (d) Lump mass by responsibility weighting:
one-sided and opposite by size.}
\label{fig:ginniemix}
\end{figure}

\begin{figure}[t]
\centering
\includegraphics[width=0.98\textwidth]{sigma_diagnosis_histograms.png}
\caption{Diagnosis under controlled kurtosis $68$. (a) The effect law:
quiet core, rare monsters. (b) Three scale estimators tell three
stories about one distribution. (c) Posterior estimates keep the tails
under the mixture prior and crush them under the Gaussian prior. (d) EM
responsibilities separate true lump pools from ordinary ones.}
\label{fig:sigmadiag}
\end{figure}

\begin{figure}[t]
\centering
\includegraphics[width=0.92\textwidth]{posterior_kurtosis_transform.png}
\caption{The posterior kurtosis transform at kurtosis $68$, $n = 150$:
tails survive estimation; lumps ride the common linear shrinkage slope
and are understated by $1 - \lambda$.}
\label{fig:posttransform}
\end{figure}


\section{One law in production and the corrected ledger}\label{sec:onelaw}

The admission null consumes the measured mixture directly. Per-stratum
pool variance enters the null variance as
$\sum_i w_i^2\phi_i^2\sigma_i^2/X^2$; the dispersion estimate deflates
each Pearson term by its own stratum's law; and tail thickness enters
through the candidate's fourth cumulant
$\lambda_4 = \sum_i a_i^4\kappa_{4,i}\sigma_i^4/\mathrm{Var}^2$, mapped
to a Gaussian-equivalent score by a strictly monotone probability match:
the null is modeled as a two-component Gaussian scale mixture with the
candidate's variance and $\lambda_4$, and the observed $z$ is carried
through its tail and an upper-tail probit. Where $K_{\mathrm{eff}}$ is
large the corrections vanish automatically; where a lump-capable pool
dominates, the deep threshold rises exactly there. A permanent
equivalence guard (one-law with a flat table reproduces the flat-gate
run exactly) protects the machinery, and the standard falsification
suite passes: the kurtosis-dialed cosine benchmarks at baseline power,
the certified counts below. Two implementation errors found and fixed
en route (a non-monotone quantile expansion; a dispersion normalizer
that halved $\hat\phi$ under the held-out split) are documented in the
lab notes; the second affected all earlier split-run numbers, so the
ledger is re-baselined here.

Currency check (this revision). All ledger pins were re-verified on the
current binary, which carries every flag-gated mechanism of
Section~\ref{sec:boundary} in its default-off state: composite
$q{=}0.30$ held-out deviance 10.488 (bit-exact, with the recorded
\texttt{DMESH\_GATE\_EXTRA\_VAR=0.35}); one-law 1527 vertices / 9.399
and one-law with law weights 2021 / 10.261 (both bit-exact); certified
counts v3${}=2$, v4${}=7$ re-run \emph{with the quadratic prior
enabled}, i.e.\ under the production configuration, unchanged. Two
caveats now on record. First, guard values are pinned per build: the
flag-off deviance of the quadratic-prior configuration showed
$\pm 0.1\%$ sensitivity to floating-point contraction across rebuilds
(10.798--10.808); vertex counts were stable throughout, so counts are
the cross-build invariant and deviances are the per-build one. Second,
the offset was certified \emph{end to end}: the production offset
recipe (plain two-exponential age ramp, coefficients fitted on the
train half) was estimated on each null dataset exactly as it would be
in production, and the full pipeline run on the result. Certified
counts: v3${}=2$ (unchanged), v4${}=3$ (quieter than the zero-offset
7): the offset absorbs level that the mesh otherwise builds, and the
gate stays honest on what remains. The audit also caught a defect
worth recording: a WAC-interaction variant of the offset, unregularized,
reached $\pm14$ logits at sparse band corners of the null datasets. It
was already dominated on real data by the plain recipe (held-out NLL
2.9365 vs 2.9408) and is excluded from production; offset designs with
interactions require regularization before use.

Ledger, post-fix. Certified-false deep faces: composite gate $2$ (v3)
and $7$ (v4), unchanged. Default gate $969$ (v3) and $996$ (v4)
on current code. Real April 2026 at $q=0.30$: composite $1857$
vertices, held-out deviance $10.488$; one-law $1527$ vertices,
$9.399$: honest pricing gives the smaller mesh and the better held-out
fit. Under the corrected law the $q$ dial recalibrates
($q=0.10$ now sits where roughly $q=0.02$ sat) and the held-out
deviance spread across the dial collapses.

Residual micro-structure the mesh cannot represent (measured; remedy
queued): a $48$ centi-logit sawtooth in residuals by WAC position
within the $0.5\%$ coupon grid, discontinuous at the Ginnie II
assignment boundary (servicing-spread position within coupon), plus a
$12$ centi-logit annual origination-seasonality cycle. Both are
covariates, not surface: coupon re-parsed from the file, spread and
seasonality into the observation offset, then the mixture law
re-estimated, since part of the fitted pool variance is this structure
in costume.


\subsection*{The completed system: joint fixed point and stratum means}

Status: measured, end of session. The law now feeds all four consumers:
the gate (per-stratum variance and tails), the IRLS working weights
(per point, including stage zero, which retires binomial-weight mega
domination at its source), the per-pool posteriors (mixture
responsibilities and effects), and the held-out metric, evaluated by
reference with the same law object rather than a refit. Surface and law
are alternated to mutual consistency with damped updates; the law
converges in about two iterations to a few-percent band, and the
residual run-to-run wobble at fixed law ($\pm10\%$ in mesh size)
lower-bounds any fixed point's tightness and is the multi-seed pass's
measured target. One instability found and designed around, one
sentence's worth: per-pool responsibility-weighted \emph{working}
weights are self-exciting (down-weight, surface departs, lumpier,
down-weight) and oscillate; weights therefore consume the stratum law,
per-pool variances are reserved for the gate. Law-priced weights
deliver held-out mixture NLL $2.977$ to $2.994$, at that point the best
predictive scores in the program, from the certification-grade gate.

A held-out Q-Q verification of the law (probit of the predictive PIT
per stratum) shows the two-component mixture near-exact through two
sigma and cutting the three-sigma tail excess roughly threefold versus
a variance-matched Gaussian, with a quantified residual ($1.3$ to
$1.9\times$) in the populated strata.

The residual histograms then exposed differing \emph{means} across size
strata, and extending the law with a per-stratum mean $\mu$ measures
the formation size gradient directly: $\mu = -0.34 / -0.02 / +0.25 /
+0.11$ across the four strata, matching the count effect in magnitude.
Consequences: the small-stratum law deflates (total $0.71\to0.60$,
lump weight $0.27\to0.21$; the $512$--$4096$ lump weight halves), the
small-pool lump one-sidedness largely dissolves (approximately
symmetric about the stratum mean; the mega fast-skew survives), and
held-out mixture NLL improves $3.000\to2.945$, the largest single
metric gain in the program, from four numbers
(Figure~\ref{fig:jfpmu}). The mega $\mu$ is unstable across the
even/odd split (a handful of giant pools dominates either half) and
wants shrinkage or cross-vintage estimation. Design conclusion: $\mu$
is signal, not noise; it belongs in a per-observation offset, one hook
serving the three measured covariates the smooth surface cannot
represent (size gradient, coupon-spread sawtooth, origination
seasonality), after which the law is re-estimated on corrected
residuals.

\begin{figure}[t]
\centering
\includegraphics[width=0.98\textwidth]{jfp_mu_diagnostics.png}
\caption{The converged one-law system with per-stratum means, held-out:
residual histograms with the mean-extended law (solid) against the
mean-zero law (dashed), and predictive Q-Q recentering on the diagonal.
Held-out mixture NLL $3.000 \to 2.945$.}
\label{fig:jfpmu}
\end{figure}

Working set: default gate (predictive bound), composite $q=0.10$
(certification) and $q=0.30$ (production), the one-law gate as
successor under validation, the permutation gate parked pending exact
permutation moments; mixture-EM parameter report, held-out mixture NLL
(headline), deviance (secondary), certified-false counts, the confound
and chirp falsifiers, and the PIT audit. Retired branches live in the
lab notes, not here.

\section{Dynamics on frozen nodes (summary)}\label{sec:dynamics}

For completeness: after the static fit, freeze the geometry and refit heights monthly
by the exact WLS normal equations (iterative sweeps leave smooth-mode residue an order
of magnitude above the theoretical noise; the monthly system is small, solve it), with
per-month reinitialization so measurement errors are independent across months. The
population kinematics are known (aging $\dot a=1$, common incentive drift), so the
transport term of the McKendrick--von Foerster description is imposed, not estimated:
build the panel in the Eulerian frame (fixed $(a,c)$ cells, jointly binned so
composition drift does not masquerade as decay) and the Lagrangian frame (fixed vintage
labels); the attached frame is the one in which the panel is separable
(cell profile $\times$ common decay $=$ rank-1 up to noise), and the decay $\lambda$ is
read off by an errors-in-variables-robust AR regression with lagged instruments in the
winning frame. This is the frame-aware upgrade of the fixed-partition autocorrelation
study: persistence quantifies signal decay (hence horizon and turnover), and the frame
answers whether the mispricing travels with the collateral or with the region of state
space.

\section{The boundary failure: measurement, forensics, and the limits of working-scale adaptivity}\label{sec:boundary}

This section reports a program of experiments on the WALA~$\approx 0$
boundary of the Ginnie surface. It corrects a measurement artifact in
our own earlier diagnostics, prices ten candidate mechanisms against
the failure, and ends at a structural conclusion about the estimator
itself. We report it at length because the failure mode is generic to
adaptive methods scored on local quadratic approximations, and because
the forensic chain, not any single fix, is the finding.

\subsection{The phenomenon, measured correctly}

Newly issued pools terminate at very low rates. In the WAC 5.75--6.25
band, the pooled termination rate at WALA 0--2 corresponds to
$-7.49$ log-odds. Every mesh configuration fitted the base of this
cliff between $-2.5$ and $-3.5$: an error of about four logits,
confined to the first eight months of pool life.

Status: our earlier diagnostics understated this error, and the
correction is itself a result. Band-averaged empirical logits with the
$+\tfrac12$ correction floor each $k=0$ pool at $-\log(2n)$; at WALA
0--2, 97\% of pools have $k=0$, so the ``observed'' dots read $-4.49$
where the pooled rate is $-7.49$. Fitted surfaces that dove past
$-4.5$ appeared to overshoot while actually improving. The proper
mixture NLL, which uses exact binomial probabilities, disagreed with
the dots throughout; the NLL was right. Reporting standard amended:
at $k=0$-saturated edges, profile dots must be pooled-rate
($\sum k/\sum n$), never averaged empirical logits.

\subsection{Two working fixes and a benchmark}

Two additions materially improved the edge. First, a locally quadratic
shrink target (Dubuc--Deslauriers along the parent edge, one-sided at
boundaries) moved the base from $-2.54$ to $-3.45$ and improved
held-out mixture NLL from 3.0000 to 2.9838 with certification
unchanged. Second, a per-pool offset column carrying a two-exponential
age ramp fitted on the training half moved the base to $-6.57$ and the
NLL to 2.9365, the best mesh result in the program; the ramp-region
NLL (2.211) essentially matches a law-weighted tensor-spline GAM
benchmark (2.159), which reaches $-8.05$ at the base by running full
IRLS on a fixed basis with an effectively unpenalized fit. The GAM
wins overall NLL (2.9045) by a thin margin distributed over the
mid-band, where it also produces 11.0 centi-logits rms of spurious
structure inside the certified window against the mesh's 4.4: the
familiar trade, prediction against policed attribution.

\begin{figure}[t]
\centering
\includegraphics[width=\textwidth]{ginnie_edge_approaches.png}
\caption{Every approach at the real WALA boundary, band-averaged
against held-out pools. All methods miss the base and bracket it from
opposite sides; differences vanish by WALA 10--12. (Dots here are
empirical-logit band averages; per the amended standard, the floored
region WALA $<8$ is re-measured with pooled-rate dots in the text.)}
\label{fig:edgeapproaches}
\end{figure}

\begin{figure}[t]
\centering
\includegraphics[width=\textwidth]{edge_failure_explained.png}
\caption{Why the pre-offset mesh failed. Top: mesh in the ramp region
before and after the offset (1{,}862 elements at median depth 17 in
the WALA $<8$ strip versus 507 at depth 14): detection was never the
problem. Bottom left: empirical-logit dots sit on the $k=0$ floor
line; pooled-rate dots reveal the true $-7.49$ base. Bottom right:
staged shrinkage deleted 42/21 centi-logits per vertex at the ramp
before the offset, 0.2 after.}
\label{fig:edgefailure}
\end{figure}

\begin{figure}[t]
\centering
\includegraphics[width=\textwidth]{ghost_bc_explained.png}
\caption{Implicit boundary conditions priced. The second-derivative
continuation from the lattice values the boundary stencil can reach
predicts $-4.0$ where $-7.5$ is needed: a resolution ceiling, not a
stencil-order problem. Ghost BCs doubled the shrinkage deletion while
improving admission, and the surface did not move.}
\label{fig:ghostbc}
\end{figure}

\begin{figure}[t]
\centering
\includegraphics[width=\textwidth]{quadratic_quality_map.png}
\caption{Quadratic quality everywhere. The 10\% trust region is
0.5--0.75 logits across the whole surface; the edge is distinguished
by the size of the required move (4.5 logits), not by uniquely bad
quadratics at one-logit scale.}
\label{fig:quadquality}
\end{figure}

\subsection{The forensic ledger}

Ten mechanisms were priced against the edge. Each is flag-gated in the
tree, each reproduces bit-stable baselines when off, and each verdict
below is measured, not argued.

\begin{center}
\begin{tabular}{p{4.6cm}p{3.1cm}p{5.6cm}}
\toprule
mechanism & edge result & verdict \\
\midrule
linear prior (baseline) & base $-2.54$ & under-dives \\
quadratic prior & base $-3.45$ & adopted; insufficient alone \\
log-WALA coordinates & worse everywhere & rejected \\
prerefine 48/96 & ramp worse & wrong layer \\
boundary $\tau$ scaling & synthetic yes, real no & real edge evidence biased at $p\to0$ \\
ghost BC, 1st/2nd derivative & real base $-2.9$ & continuation capped by interior resolution ($-4.0$ ceiling, measured) \\
untaxed converged sweeps & base $-2.80$ & proposals themselves stall at $-4.03$ \\
admit-all (gate open) & base $-2.93$ & allocator, not the test, starves the edge \\
honest decision scores only & base $-3.05$ & allocation moves, commits do not \\
manual age-ramp offset & base $-6.57$ & works; per-dataset crutch \\
\bottomrule
\end{tabular}
\end{center}

\subsection{Why: the trust region and the misallocated information}

The edge strip (WALA $<8$, WAC 5--7) contains \emph{more} total Fisher
information than a same-width mid-band strip (6{,}425 vs 3{,}769 units
of $\sum n_i p_i(1-p_i)$); the band-mean logit is pinned to
$\pm 0.0125$. The data is sufficient. What fails is the currency in
which every internal decision is denominated. Measured across all
120{,}595 pools at the production fit: the local quadratic model of
the pool log-likelihood has a 10\%-accuracy trust region of 0.5--0.75
logits essentially everywhere on the surface; 43\% of pools exceed
20\% error already at a one-logit move; only 4\% of pools have a
four-logit trust region. Mid-band this is harmless because converged
fits make candidate moves centi-logit sized. At the edge the required
move is 4.5 logits, nine trust radii. One $k=0$ pool with $n=150$
gains 7.6 nats moving from $-2.95$ to $-7.5$; its working-scale score
for the same move is a fraction of a nat. Refinement ranking
($\hat\delta^2 x^Tx$), gate admission, and EB commits all consume this
score, so the budget cascade, the admissions, and the committed
heights starve the edge coherently. The admit-all experiment isolates
the allocator: with the FDR gate fully open, the boundary column
became \emph{shallower}, because candidate generation and vertex caps
spend the budget by the same underpriced score before boundary
candidates exist.

The general statement: adaptive methods that commit irreversible
resources on local-quadratic scores are safe exactly where adaptation
has already brought the fit within the trust radius, which is circular
precisely where a large feature has not yet been found. The GAM
escapes not because its quadratics are better but because it makes no
irreversible decisions: its basis exists a priori and its IRLS
iterates every quadratic to convergence before anything is read off.

\subsection{Why it is hard to fix in place}

An exact tempered profile likelihood per vertex coordinate (concave;
one-dimensional Newton; temper $c_i = (1+s^2_{s(i)} n_i p_i(1-p_i))^{-1}$
reproducing the mixture-law weighting) was implemented and integrated
three ways. All three degraded the global fit (held-out NLL 3.04,
3.85, 3.31 against the 2.98 baseline). The reason is structural. The
production sweep is a hybrid: working responses are frozen at the last
IRLS refresh while neighbor heights update live (Gauss--Seidel on a
deliberately stale linearization), and an empirical-Bayes tax applies
every sweep. This hybrid has no likelihood-space equivalent. An exact
update evaluated on live state bypasses the damping the hybrid relies
on (removing the per-sweep tax alone degraded every global metric,
so the tax is load-bearing); an exact update evaluated on frozen
state fights the live Gauss--Seidel; and a trust-region trigger makes
vertices oscillate between the two semantics. Substituting honest
scores only at decision points, with commits untouched, moves
allocation (boundary admission and depth rise) but the quadratic
commits still park boundary vertices at information-weighted star
averages, and honest nat-scale gains fed to a null calibrated for the
old statistic reallocate budget toward mega-pool texture. Iterating
warm-started quadratics to convergence is exactly correct at the level
of any single one-dimensional fit (for binomials the converged working
fit \emph{is} the profile MLE), but making every decision wait for an
outer-converged state at every resolution round is
solve--estimate--mark--refine with converged solves: functionally a
new solver, implemented as scheduling discipline bolted onto machinery
designed for staleness.

\subsection{Exact-solver resolution and the remaining calibration problem}

The proposed likelihood-space resolution has now been implemented. B1 replaced stage fitting by
coordinate ascent on the exact tempered penalized objective, unified active and inactive hierarchical
columns, required every candidate commit to realize its profiled gain, and replaced nodal-star
retirement by the hierarchy-aware rule of Section~\ref{sec:lifecycle}. On the original split the
closed B1 regression gave held-out NLL $2.9335$ and a WALA $0$--$2$ fitted logit of $-6.96$ against
the law-marginal reference $-7.95$. Seeds 202 and 404 remained stable, and clean reruns were
byte-identical apart from timing text.

B0 was then repeated against the completed exact solver. The saturation-triggered exact fallback was
proved inert and deleted, together with an unused pause hook. The height guard and IRLS step cap were
not deleted: although they no longer control the stage optimizer, they still change the legacy gate's
candidate scores and topology. After the dead fallback was removed, the 20-split battery had zero
eruptions, mean NLL $2.9449$, sd $0.0054$, and worst split $2.9537$. These remaining guards are
legacy calibration dependencies, not accepted principles; they should disappear only when the gate
itself no longer consumes the working-scale path.

B2 introduced the exact signed-root gain \eqref{eq:exactgain}. It is deterministic, nonnegative in
squared magnitude, agrees with the old score in the quadratic regime, and exposes large direction and
ranking disagreements away from that regime. Shadow mode leaves the accepted B1 fit unchanged. B3
then attempted to calibrate this score through the candidate-specific sandwich
\eqref{eq:b3score}. The proposal looked good on complete-null scale diagnostics and improved the
three real-data regressions, but the predeclared v3/v4 instruments rejected it as recorded in
Section~\ref{sec:certnull}. Prediction and numerical stability have therefore transferred to the
exact solver; certification has not.

The current status is consequently asymmetric but clear:
\begin{itemize}[leftmargin=2em,itemsep=1pt]
\item exact likelihood fitting, exact candidate profiles, and hierarchical lifecycle: implemented and
regression tested;
\item eruption prevention and boundary access: passed;
\item signed-root exact gain as the candidate currency: selected and shadow validated;
\item null calibration and false-discovery certification for that currency: failed B3 and remains open;
\item certified counts elsewhere in this document: legacy-gate results only.
\end{itemize}
The next experiment is the conformalized held-out gain test of
Section~\ref{sec:certnull}, not another post-hoc rescaling of the same global score family.

\section{Discussion}\label{sec:discussion}

Open items, with status. (i) Exact-score calibration: B3's candidate-specific
sandwich is locally well behaved but fails the certified windows under the global
empirical-null prefix rule. (ii) Heavy-tail complete-null calibration: the clean binomial
null shuts completely, while a matched pool-effect law produces occasional coarse false
branches that are amplified by repeated adaptive testing. The immediate experiment is the
one-family/frozen-dictionary/full-adaptive decomposition above, followed by branchwise or
sequential error spending if the diagnosis survives a larger seed battery. (iii)
Conformalized held-out gain certification remains a separate, more expensive route and is
designed but not implemented. (iv) Issuer-clustered replication is specified and exposed
by v4, but not yet built into the exact gate. (v) The remaining height and IRLS guards are
active only through the legacy gate and must be re-tested for deletion after calibration
migrates. (vi) Pool-law identification ultimately requires loan-level disclosure or a
panel; formation covariates can narrow the bundle in the meantime. Held-out evidence in
this revision is mixed: the exact-solver stability claim uses twenty resampled splits,
while many older figures remain single-split illustrations as labeled.

The method is designed for a narrow but common problem: correcting an existing probability model on
thin, irregular support while making each added degree of freedom auditable. Its main practical
advantage is not merely adaptivity. It is the combination of geometric adaptivity with an explicit
multiplicity gate, hierarchical shrinkage, and a diagnostic layer that refuses to label all residual
variation as discoverable structure.

The strongest numerical conclusion is that the attribution boundary is resolution dependent but
interpretable. Underresolved spatial structure enters the persistent component; as information grows
or the coordinate metric improves, that variation is reclaimed by the surface. Continuous WACLS
pseudo-counts preserve this behavior for balance-weighted outcomes and avoid arbitrary integer
quantization.

Several limitations remain. Calendar-time shocks are omitted from the current generative model;
whole-run FDR control for the adaptively generated hypothesis sequence is not proved; the exact-score
gate has not passed the certified-window instruments; conformal validity would require explicit
pool/issuer exchangeability and a candidate family frozen before calibration; selection-aware
estimation of the EB variance remains unsettled; and no adaptive approximation-rate theorem is
claimed. True balance-weighted use in the Ginnie Mae application also requires current loan-level
balances, because pool-level disclosures do not contain current-balance WACLS.

These limitations suggest a staged next step rather than a change in the central design.  First isolate
root-family tail calibration from refitting and descendant generation on the heavy-tail complete null,
and test a branchwise sequential budget against the same frozen seed battery.  Then extend the working
model to calendar effects, issuer blocks, and continuous quasi-binomial exposure and study the gate and
attribution moments under that enlarged dependence structure.

\begin{thebibliography}{9}\small
\bibitem{efron2004} B.~Efron, Large-scale simultaneous hypothesis testing: the choice
of a null hypothesis, \emph{JASA} 99 (2004) 96--104.
\bibitem{efron2008} B.~Efron, Microarrays, empirical Bayes and the two-groups model,
\emph{Statist.\ Sci.} 23 (2008) 1--22.
\bibitem{bh1995} Y.~Benjamini and Y.~Hochberg, Controlling the false discovery rate,
\emph{JRSS-B} 57 (1995) 289--300.
\bibitem{abramovich2006} F.~Abramovich, Y.~Benjamini, D.~Donoho, I.~Johnstone, Adapting
to unknown sparsity by controlling the false discovery rate, \emph{Ann.\ Statist.} 34
(2006) 584--653.
\bibitem{mitchell1991} W.~F.~Mitchell, Adaptive refinement for arbitrary finite-element
spaces with hierarchical bases, \emph{J.\ Comput.\ Appl.\ Math.} 36 (1991) 65--78.
\bibitem{binev2004} P.~Binev, W.~Dahmen, R.~DeVore, Adaptive finite element methods
with convergence rates, \emph{Numer.\ Math.} 97 (2004) 219--268.
\bibitem{dl1986} R.~DerSimonian and N.~Laird, Meta-analysis in clinical trials,
\emph{Controlled Clinical Trials} 7 (1986) 177--188.
\bibitem{firth1993} D.~Firth, Bias reduction of maximum likelihood estimates,
\emph{Biometrika} 80 (1993) 27--38.
\bibitem{buhlmann1967} H.~B\"uhlmann, Experience rating and credibility,
\emph{ASTIN Bulletin} 4 (1967) 199--207.
\bibitem{henderson1975} C.~R.~Henderson, Best linear unbiased estimation and
prediction under a selection model, \emph{Biometrics} 31 (1975) 423--447.
\bibitem{paule1982} R.~C.~Paule and J.~Mandel, Consensus values and weighting factors,
\emph{J.\ Res.\ Natl.\ Bur.\ Stand.} 87 (1982) 377--385.
\bibitem{efron2004pred} B.~Efron, The estimation of prediction error: covariance
penalties and cross-validation, \emph{JASA} 99 (2004) 619--632.
\bibitem{ye1998} J.~Ye, On measuring and correcting the effects of data mining and
model selection, \emph{JASA} 93 (1998) 120--131.
\bibitem{vaupel1979} J.~Vaupel, K.~Manton, E.~Stallard, The impact of heterogeneity in
individual frailty on the dynamics of mortality, \emph{Demography} 16 (1979) 439--454.
\bibitem{vovk2005} V.~Vovk, A.~Gammerman, and G.~Shafer, \emph{Algorithmic Learning in a Random
World}, Springer, 2005.
\end{thebibliography}

\end{document}
